\documentclass{article}

\usepackage[left=30mm,right=10mm,top=20mm,bottom=20mm,papersize={210mm,297mm}]{geometry}

% межстрочный интервал
\linespread{1.5}

% кликабельные референсы
\usepackage[hidelinks]{hyperref}

% абзацный отступ
\usepackage{indentfirst}

% top page pagination
\usepackage{fancyhdr}
\fancyhf{}
\pagestyle{fancy}
\fancyfoot[C]{\thepage}
\renewcommand\headrulewidth{0pt}

% times new roman support
\usepackage{fontspec}
\setmainfont{Times New Roman}

\usepackage{polyglossia}
\setdefaultlanguage{russian}
\setotherlanguage{english}
% Настройка моноширинного шрифта (для команд типа \texttt)
\setmonofont[Scale=MatchLowercase]{Consolas}
\newfontfamily\cyrillicfonttt[Scale=MatchLowercase]{Consolas}

% more font sizes
\usepackage{scrextend}

% Formatting TOC
\usepackage{tocloft}
\renewcommand{\cftsecfont}{\normalfont\mdseries}% titles in non-bold
\renewcommand{\cftsecpagefont}{\normalfont\mdseries}% page numbers in non-bold
\renewcommand{\cftsecleader}{\mdseries\cftdotfill{\cftdotsep}}% dot leaders in non-bold

% --- ПОЛНОСТЬЮ ОТКЛЮЧАЕМ АВТОМАТИЧЕСКИЙ ЗАГОЛОВОК ---
\makeatletter
\renewcommand{\@cftmaketoctitle}{}
\makeatother
% ----------------------------------------------------

\renewcommand{\cftsecaftersnum}{.}
\renewcommand{\cftsubsecaftersnum}{.}

\usepackage{titlesec}
\titleformat{\section}{\normalfont\center\bfseries\uppercase}{}{0em}{}
\titleformat{\subsection}{\normalfont\center\bfseries}{\thesubsection}{0.5em}{}
% \titlespacing{\section}{0pt}{0.2\baselineskip}{0.2\baselineskip}
% \titlespacing{\subsection}{0pt}{0.2\baselineskip}{0.2\baselineskip}

% remove footnote margin
\usepackage{footmisc}
\usepackage{url}
\urlstyle{same}

% tables support
\usepackage{float}
\usepackage{multirow}
\usepackage{tabularx}
\usepackage{tablefootnote}
\usepackage{caption} 
% custom captions
\renewcommand{\figurename}{Рис.}
\renewcommand{\tablename}{Таблица}
\captionsetup[table]{labelsep=endash,position=top,justification=raggedright,singlelinecheck=false,font=normalsize}
\captionsetup[figure]{labelsep=endash,justification=centering,singlelinecheck=false,font=normalsize}

% image support
\usepackage{graphicx}
\graphicspath{ {images/} }

% math symbols (R, C, N)
\usepackage{amssymb}


% plots
\usepackage{tikz}
\usetikzlibrary{shapes.geometric, arrows.meta, positioning, fit, backgrounds, calc}
\usepackage{pgfplots}
\usepackage{subcaption}
\usepackage{graphicx}
\pgfplotsset{compat=1.17}


% disable hyphenation
\tolerance=1
\emergencystretch=\maxdimen
\hyphenpenalty=10000
\hbadness=10000

% prevent footnote spreading onto next page
\interfootnotelinepenalty=10000

% не курсивных текст в формулах через \text{}
\usepackage{amsmath}

% [1] нумерация в списке литературы
\usepackage{enumitem}

\usepackage{threeparttable}

% Пакет для работы с библиографией по ГОСТу
\usepackage[
    backend=biber,
    bibstyle=gost-numeric,
    citestyle=numeric-comp,
    language=auto,
    autolang=other,
    sorting=none,
    defernumbers=true
]{biblatex}

% 1. Меняем нумерацию с "1." на "[1]"
\DeclareFieldFormat{labelnumberwidth}{\mkbibbrackets{#1}}

% 2. Меняем "URL" на "Режим доступа" для русского языка
\DefineBibliographyStrings{russian}{
  url = {Режим доступа\addcolon},
  urlseen = {дата обращения}
}
\DeclareFieldFormat{url}{Режим доступа\addcolon\space\url{#1}}
\DeclareFieldFormat{urldate}{\mkbibparens{дата обращения\addcolon\space#1}}
% Подключаем файл с вашей базой литературы (создайте его рядом с .tex файлом)
\addbibresource{references.bib}

\usepackage{setspace}

% =============================================


\title{CourseWork}
\author{}
\date{}

\begin{document}

% абзацный отступ
\setlength{\parindent}{1.25cm}
% remove footnote margin
\setlength{\footnotemargin}{5pt}
% text justification
\sloppy
\frenchspacing % fixes double spaces after commas

\thispagestyle{empty}
\begin{center}
    \KOMAoption{fontsize}{13pt}

    ФЕДЕРАЛЬНОЕ  ГОСУДАРСТВЕННОЕ АВТОНОМНОЕ ОБРАЗОВАТЕЛЬНОЕ
    УЧРЕЖДЕНИЕ \\
    ВЫСШЕГО ПРОФЕССИОНАЛЬНОГО ОБРАЗОВАНИЯ \\
    «НАЦИОНАЛЬНЫЙ ИССЛЕДОВАТЕЛЬСКИЙ УНИВЕРСИТЕТ \\
    «ВЫСШАЯ ШКОЛА ЭКОНОМИКИ»

    \bigbreak
    \bigbreak

    \textit{Факультет информатики, математики и компьютерных наук} \\
    % \textbf{Программа подготовки бакалавров по направлению} \\
    % \textbf{01.03.02 Прикладная математика и информатика} \\

    \bigbreak
    \bigbreak

    \textit{Тюрин Владислав Николаевич} \\

    \KOMAoption{fontsize}{14pt}
    \bigbreak
    \textbf{Построение графа знаний на основе пользовательских данных с возможностью экспорта в Obsidian} \\
    \bigbreak
    Выпускная квалификационная работа - БАКАЛАВРСКАЯ РАБОТА \\
    по направлению подготовки 01.03.02 Прикладная Математика и Информатика \\
    образовательная программа \\
    «Прикладная Математика и Информатика» \\

    \KOMAoption{fontsize}{16pt}
    
\end{center}

\KOMAoption{fontsize}{14pt} % font size = 14
\renewcommand{\footnotesize}{\fontsize{14pt}{14pt}\selectfont} % footnote font size = 14
\vspace{50pt}

\noindent
\begin{minipage}[t]{0.45\textwidth}
    Рецензент \\
    \\ % Пустая строка для выравнивания, если нужно место под ФИО
    % \underline{\hspace{5cm}}
\end{minipage}
\hfill
\begin{minipage}[t]{0.5\textwidth}
    \begin{flushright}
        Научный руководитель \\
        к. ф. – м. н., старший преподаватель \\
        \underline{\hspace{5cm}} \\
        В. М. Кочеганов
    \end{flushright}
\end{minipage}

\begin{center}
    \vspace*{\fill}
    Нижний Новгород, 2026
\end{center}


\clearpage

\begin{center}
    \textbf{ОГЛАВЛЕНИЕ}
\end{center}

\tocloftpagestyle{fancy}
\tableofcontents


\clearpage

\clearpage

\begin{center}
    \section*{Аннотация}
    % \addcontentsline{toc}{section}{Аннотация}
\end{center}

\begin{center}
    \section*{Abstract}
    % \addcontentsline{toc}{section}{Abstract}
\end{center}

\begin{center}
    \section*{Введение}
    \addcontentsline{toc}{section}{Введение}
\end{center}


В эпоху информационной перегрузки
возможность эффективно управлять данными
стала особенно важной.
Классических способов работы с информацией,
таких как ведение заметок,
уже недостаточно,
особенно для специалистов
регулярно работающими с огромными массивами многотемных
данных. В качестве решения были разработаны специализированные системы
для управления знаниями, например, Obsidian.
Они берут на себя роль «второго мозга»,
освобождая пользователя от хранения всей информации у себя в голове
и упрощая манипуляции с ней.

Современные системы управления знаниями поддерживают
генерацию с дополненной выборкой (Retrieval-Augmented Generation, RAG).
Большинство RAG-систем используют гибридный поиск
(комбинацию лексического и семантического).
Такой подход имеет несколько слабостей:
лексический поиск не умеет искать синонимы,
а семантический может упустить релевантные данные из-за большой разницы
в векторных представлениях (эмбеддингах).

Более продвинутой стратегией поиска является
использование графов знаний (ГЗ),
которые соединяют семантически связанные сущности.
Каждая связь имеет направление и тип,
отражающий семантическую связь.
Это открывает доступ к многошаговому поиску,
что решает проблемы гибридного поиска и
позволяет увеличить полноту результатов.

Несмотря на преимущества ГЗ,
их построение и поддержка требуют колоссальных
ресурсов. По мере роста базы знаний возникает «узкое горлышко»:
пользователям недостаточно когнитивных ресурсов и времени
на анализ и связывание заметок.

На данный момент существующие решения
(графовые базы данных корпоративного уровня или
облачные приложения )
имеют критические недостатки.
Базы данных требуют дорогостоящей инфраструктуры
и чрезмерно зависят от больших языковых моделей (LLM),
что приводит к генерации бессмысленных связей и галлюцинациям.
Облачные решения, несмотря на их удобство для повседневного использования,
создают риск утечки данных и принудительно привязывают пользователей к поставщику.

Цель работы -- разработка автоматизированного алгоритма
для построения графов знаний путем семантического связывания заметок в экосистеме Obsidian,
а также сравнение разработанного решения с существующими аналогами в сценариях использования RAG.

Для достижения поставленной цели были сформулированы следующие задачи:

\begin{enumerate}
    \item Провести анализ существующих решений для построения графов знаний
          для выявления их ограничений и функциональных особенностей.
    \item Разработать алгоритм автоматического поиска и типизации связей, включающий:
          извлечение содержимого Markdown-файлов; гибридный поиск кандидатов; классификацию связей; сохранение
          найденных отношений, поддерживающее инкрементальные обновления.
    \item Создать систему валидации и репрезентативный аннотированный набор данных для объективной оценки качества работы алгоритма
          и его сравнения с аналогами.
    \item Провести итеративное улучшение решения путем экспериментов с различными моделями,
          гиперпараметрами и методами фильтрации связей
          для повышения как точности системы, так и ее вычислительной эффективности.
\end{enumerate}

Объект исследования — процессы поиска,
структурирования и извлечения информации в системах управления знаниями.
Предмет исследования — методы и алгоритмы автоматического
семантического связывания текстовых данных для построения ГЗ.

Алгоритм, полученный в результате данной работы,
ориентирован на продвинутых пользователей баз знаний,
которым может быть важна суверенность данных,
возможность работы алгоритма локально
или при ограниченных вычислительных мощностях,
а также эффективно обновлять ГЗ.
Это подтверждает практическую значимость работы.

% МОЙ ТЕКСТ
% Введение
% (Пишется обычно в самом конце, но каркас нужен сразу)
%     Актуальность темы: Информационная перегрузка, рост PKM-систем, ограничения гибридного поиска в RAG и сложность ручного ведения Knowledge Graphs (KG).
%     Цель работы: Разработка автоматизированного пайплайна для извлечения, фильтрации и типизации связей между заметками в экосистеме Obsidian.
%     Задачи: (Прямо из твоего пропоузала: анализ решений, разработка алгоритма, создание датасета и фреймворка оценки, тестирование).
%     Объект и предмет исследования.

% Глава 1. Теоретические основы и обзор предметной области (Теория)
% Здесь мы раскрываем твои пункты 1, 2, 3 и 4. Глава должна плавно подвести читателя к тому, что существующие решения не идеальны, и твой подход необходим.
%     1.1. Управление персональными знаниями (PKM) и графы знаний
%         Концепция "Второго мозга" и экосистема Obsidian.
%         Роль Knowledge Graphs (KG) в структурировании информации.
%         Проблема "бутылочного горлышка" при ручном ведении масштабных PKG (Personal Knowledge Graphs).
%     1.2. Технологии поиска и обработки естественного языка в RAG
%         Ограничения классического лексического (BM25) и семантического (Cosine Similarity) поиска. Концепция гибридного поиска.
%         Применение Retrieval-Augmented Generation (RAG) в системах управления знаниями.
%         Обзор языковых моделей: архитектура трансформеров, LLM (напр., Gemini, Llama) и легковесные NLI-модели (Cross-encoders, DeBERTa).
%     1.3. Анализ существующих решений для построения графов знаний
%         Enterprise-решения: Neo4j, Memgraph (плюсы: масштаб, минусы: статика, оверхед, галлюцинации LLM).
%         Академические пайплайны: EduKG (плюсы: микро-графы, эффективность, минусы: привязка к онтологиям).
%         Облачные PKM и плагины: Mem.ai, Smart Connections, Khoj (плюсы: интеграция, минусы: vendor lock-in, "черный ящик", отсутствие сложной семантики).
%     1.4. Методы оценки качества графов и систем RAG
%         Метрики классификации для проверки качества связей (Precision, Recall, F1-score).
%         Фреймворки оценки LLM и RAG-систем (в частности, RAGAS: Context Precision, Answer Faithfulness).



% Глава 3. Экспериментальное исследование и оценка результатов (Практика: Часть 2)
% Здесь ты доказываешь, что твой код работает лучше аналогов.
%     3.1. Дизайн эксперимента и Baseline-решение
%         Описание тестового стенда и гиперпараметров.
%         Создание базового решения (baseline) на основе простого dense retrieval + LLM для сравнения.
%     3.2. Оценка качества извлечения и типизации связей
%         Сравнение пайплайна с baseline по метрикам Precision и Recall на собранном датасете.
%         Анализ ошибок (где алгоритм ошибся и почему).
%     3.3. Эффективность построенного графа в задачах RAG
%         Оценка влияния сгенерированных графов на качество ответов RAG.
%         Использование RAGAS (Context Precision, Answer Faithfulness).
%     3.4. Оценка вычислительной эффективности
%         Анализ затрат (скорость работы локальных моделей vs. стоимость/скорость API).
%         Ограничения разработанного решения.

% Заключение
%     Краткие выводы по каждой главе.
%     Подтверждение того, что цель достигнута (пайплайн создан, когнитивная нагрузка снижена, RAG улучшен).
%     Направления для дальнейших исследований (Future work).

% Список литературы
% (Сюда отлично лягут все твои источники: Lewis, Gao, Edge, Nogueira, He и документации к Neo4j/LanceDB).


% % Объект работы
% Объект работы --- мультимодальное машинное обучение.

% % предмет работы
% Предмет работы --- применение мультимодального машинного обучения к задаче
% классификации вида больной/здоровый человек.

% % Цель работы
% Цель работы --- изучение применимости мультимодального машинного обучения на основе структурированных
% и неструктурированных данных к задаче классификации больных и здоровых людей.

% % Задачи
% \begin{itemize}
%     \item[] Задачи:
%           \begin{itemize}
%               \item[$\bullet$] Изучить текущий прогресс машинного обучения в области медицины
%               \item[$\bullet$] Изучить процесс разработки мультимодальной модели
%               \item[$\bullet$] Изучить виды слияния модальностей
%               \item[$\bullet$] Найти мультимодальный датасет, состоящий из
%                     структурированных и неструктурированных данных
%               \item[$\bullet$] Изучить имеющиеся решения на выбранном датасете
%               \item[$\bullet$] Выбрать решение, предложить к нему улучшения,
%                     провести над ним эксперименты
%               \item[$\bullet$] Проанализировать результаты, сделать выводы
%           \end{itemize}
% \end{itemize}


\clearpage
\section{Теоретические основы и обзор предметной области}
\subsection{Системы управления персональными знаниями}

В эру информационной перегрузки
необходимо эффективно управлять данными.
Традиционные методы хранения и обработки информации,
такие как рукописные заметки или их цифровые аналоги,
хороши для повседневного использования, но
не способны обеспечить необходимую функциональность
для специалистов, работающих с огромными массивами информации.
Возрастающая актуальность этой проблемы
стимулировала развитие специальных систем для управления персональными знаниями
(Personal Knowledge Management, PKM).


В основу PKM-систем легли результаты множества исследований
из области информационной архитектуры и когнитивной психологии.
В последнее время большое влияние оказывает
концепция «Второго мозга» (Second Brain),
популяризированная исследователем Тиаго Форте~\cite{Forte_2022}.
Эта концепция позволяет переложить хранение и организацию информации на внешние носители,
тем самым разгружая рабочую память человека.

Человеческое мышление ассоциативно,
в нем одна концепция может принадлежать разным областям знания.
Иерархические структуры не обладают достаточной гибкостью,
так как требуют однозначной категоризации,
что противоречит устройству человеческого мышления.
Их все чаще заменяют на
сетевые структуры, наподобие
Zettelkasten («картотечный ящик»)~\cite{Ahrens_2017}.
Этот подход, разработанный социологом Никласом Луманом,
заключается в хранении небольших атомарных заметок,
связанных друг с другом.
Связи в нем представляют основную ценность:
информация, хранящаяся в заметках, не абстрагирована,
а встроена в общий пласт знаний,
что ускоряет поиск и позволяет делать сложные,
на первый взгляд, неочевидные умозаключения.

Популярным представителем современной реализации подхода Zettelkasten
является программа Obsidian~\cite{Obsidian}.
В ней заметки хранятся в виде файлов формата Markdown,
которые пользователь может разместить в удобном ему месте.
Заметки поддерживают двунаправленное связывание путем
упоминания заметки в тексте (формат [[wikilink]]).
Кроме того, в заметках можно добавлять специальные пометки, теги,
по которым можно осуществлять быструю фильтрацию и поиск.
Такая архитектура облегчает организацию данных,
а также обеспечивает их суверенитет, поскольку
пользователь сам выбирает место их хранения.

В основе сетевых PKM-систем лежит концепция персонального графа знаний.
Граф знаний представляет собой структуру,
объединяющую информацию о сущностях и
взаимосвязи между ними.
Его узлы~--- это сущности, например, объекты, люди, термины, идеи и так далее.
Ребра~--- направленные типизированные связи.

\begin{figure}[H]
    \centering
    % Масштабируем рисунок точно по ширине текста
    \resizebox{\textwidth}{!}{%
        \begin{tikzpicture}[
                entity/.style={circle, draw, thick, minimum size=2.5cm, align=center, font=\bfseries},
                literal/.style={rectangle, rounded corners, draw, thick, minimum height=1cm, align=center},
                edge/.style={->, >=Stealth, thick},
                label/.style={fill=white, font=\scriptsize, inner sep=1.5pt}
            ]

            % --- УЗЛЫ (Nodes) ---
            % Синий кластер
            \node[entity, fill=blue!10, draw=blue!80] (bob) {Боб};
            \node[entity, draw=blue!80, above left=2.5cm and 2.5cm of bob] (alice) {Алиса};
            \node[literal, draw=blue!80, below left=3cm and 1.5cm of bob] (person) {Человек};
            \node[literal, draw=blue!80, below right=3cm and 1.5cm of bob] (date) {14 июля 1990};

            % Красный кластер (чуть-чуть сдвинули обратно, чтобы сэкономить ширину)
            \node[entity, fill=red!10, draw=red!80, right=4cm of bob] (monalisa) {Мона Лиза};

            % Зеленый кластер
            \node[entity, draw=green!60!black, above right=2.5cm and 1.5cm of monalisa] (davinci) {Леонардо \\ да Винчи};
            \node[literal, draw=green!60!black, below right=2.5cm and 1.5cm of monalisa] (video) {Видео о \\ Джоконде};

            % --- РЕБРА (Edges) ---
            \draw[edge, blue!80] (bob) -- node[label, sloped] {друг} (alice);
            \draw[edge, blue!80] (bob) -- node[label, sloped] {интересуется} (monalisa);
            \draw[edge, blue!80] (bob) -- node[label, sloped] {экземпляр класса} (person);
            \draw[edge, blue!80] (bob) -- node[label, sloped] {дата рождения} (date);

            \draw[edge, green!60!black] (monalisa) -- node[label, sloped] {автор} (davinci);
            \draw[edge, green!60!black] (video) -- node[label, sloped] {предмет} (monalisa);

        \end{tikzpicture}%
    }
    \caption{Пример графа знаний}
    \label{fig:KG_example}
\end{figure}

Подобная структура данных позволяет создавать
кластеры (сообщества) из семантически близких сущностей,
осуществлять поиск на основе связей.
Это имеет особую ценность в системах управления знаниями, так как
превращает разрозненные пользовательские заметки в
связные группы, что делает поиск и работу с данными эффективнее и результативнее.

\subsubsection*{Простые и семантически типизированные связи}

Как правило, в ГЗ все ребра имеют направление и тип,
который отражает характер отношения между сущностями.
Это обеспечивает более продвинутую навигацию, позволяя
отсеивать нерелевантные направления на раннем этапе,
что открывает возможность использовать многошаговые запросы.

В механизме wikilink-ссылок, который использует Obsidian,
по умолчанию нет поддержки типизации.
Однако она может быть добавлена с помощью популярного расширения
DataView~\cite{dataview_plugin}, которое позволяет задавать тип связи при помощи следующего синтаксиса:
\texttt{link\_type:: [[wikilink]]}.

\subsubsection*{Специфика целевой аудитории}

Основной аудиторией систем управления знаниями
являются люди, чья деятельность
тесно сопряжена с обработкой больших массивов данных.
Исследователи, журналисты и другие специалисты
используют PKM системы для агрегации данных, создавая краткие выжимки,
для их трансформации, находя пробелы в знаниях или данных,
и множества других целей.
Например, эти системы могут быть использованы студентами
для подготовки ответов на экзаменационные вопросы на основе материалов занятий.

\subsubsection*{Постановка проблемы: Когнитивное «узкое место»}

У ГЗ есть недостаток, заключающийся в ручном процессе извлечения и связывания информации.
На начальных этапах создание ссылок и их типизация
происходит естественным образом и не требует особых усилий.
Однако при преодолении определенного порога по объему информации,
например, при наличии в хранилище сотен или тысяч заметок
пользователь утрачивает возможность удержать весь контекст в голове.
Из-за этого процесс связывания становится трудозатратным и неточным:
пользователь начинает пропускать связи, особенно неочевидные.
Как следствие, граф становится разреженным и деградирует качество поиска, так как
из-за отсутствующих связей алгоритму не удается найти необходимую информацию.
В итоге граф начинает терять свою ценность.

Таким образом, алгоритм, способный автоматически
анализировать семантику заметок, находить связи между ними
и типизировать их, тем самым поддерживая актуальность и плотность
графа знаний, устранит это «узкое место»
и освободит когнитивные ресурсы пользователей под другие задачи.

\subsection{Технологии поиска и обработки естественного языка в системах RAG}

Использование генеративных моделей в чистом виде
влечет высокие риски «конфабуляций» и «галлюцинаций»:
ответ модели может казаться правдоподобным, но содержать фактологические
или причинно-следственные ошибки~\cite{Ji_2023}.

Для решения этой проблемы была разработана технология
Retrieval-Augmented Generation (RAG) \cite{Lewis_2020}.
Она добавляет дополнительный этап поиска информации перед генерацией ответа.
На основе запроса пользователя алгоритм извлекает наиболее релевантную
информацию из доступных источников, таких как текстовые файлы или базы данных,
после чего передает ее языковой модели в качестве контекста.
Такой подход заставляет LLM выступать в роли анализатора,
а не источника фактов.

Эффективность RAG-систем зависит от качества алгоритмов извлечения.
В связи с этим, существует несколько подходов к поиску необходимых данных.
Наиболее популярными методами являются лексический и семантический поиск.

Лексический поиск основан на точных совпадениях
лексем из запроса пользователя с содержащимися в хранилищах.
Данный подход отлично справляется с поиском специфических терминов,
имен собственных и прочих точных формулировок, но
не справляется, например, в случае замены слов на синонимы или чередующихся корней.
В случае нахождения лексемы в корпусе из него извлекается
фрагмент из некоторого количества предшествующих и последующих слов или предложений.
Популярным алгоритмом лексического поиска является BM25~\cite{Robertson_2009}.

Для преодоления этих ограничений используется семантический поиск,
основанный на векторном сходстве~\cite{Karpukhin_2020}.
Исходный текст разбивается на фрагменты (чанки), для каждого из которых вычисляется
многомерный числовой вектор (эмбеддинг) из одного пространства.
Поиск происходит путем вычисления сходства между этими векторами
с помощью метрик, таких как косинусное сходство.
Если полученное значение больше некоторого порога, то
чанки считаются релевантными, и текст, соответствующих эмбеддингу,
извлекается.
Этот подход позволяет улавливать смысл текста,
не основываясь на конкретных словах или формулировках.

Векторный поиск также имеет несколько недостатков~\cite{Gao_2023}:
\begin{itemize}
      \item Потеря конкретики: смысловое значение фрагментов усредняется,
            что ведет к потере деталей, таких как имен собственных, использованных терминов, и так далее.
            Чем больше чанк, тем больше деталей теряется.
      \item Проблема «дальних эмбеддингов»: схожие данные могут иметь сильно
            отличающиеся друг от друга векторы, из-за чего могут не быть найдены.
\end{itemize}

Современные RAG системы комбинируют эти подходы в
так называемый гибридный поиск~\cite{Gao_2023}. Он позволяет
компенсировать недостатки одного подхода с помощью другого,
что потенциально может увеличить полноту поиска.
В нем извлеченные чанки получают унифицированную оценку
и ранжируются по ней.

Несмотря на все преимущества этих подходов,
обработка сложных аналитических запросов,
требующих многошагового рассуждения,
требует более сложных алгоритмов поиска.
В случаях, когда нужная информация напрямую не связана с запросом,
она не может быть найдена
ни семантическим, ни лексическим поиском, из-за чего
поисковый алгоритм RAG не сможет извлечь необходимые данные.

Более продвинутой стратегией поиска является Graph RAG~\cite{Edge_2024}.
В ней в алгоритм поиска дополняется обходом графа знаний.
Для первоначального сбора вершин используется гибридный поиск.
Затем начинается итеративный обход соседей на основе
имеющихся ребер и их типов, что позволяет системе собирать дополнительный
контекст из файлов, которые могут не содержать ключевых слов, или быть семантически далекими,
но при этом логически связанными и обладающими необходимой информацией,
что позволяет значительно расширить контекст.

\textbf{Языковые модели в задаче связывания.}
Поскольку необходимо искать тексты, обладающие не только высоким семантическим сходством
или содержащие ключевые слова, нужно использовать иные от сравнения эмбеддингов подходы для анализа
связности текстов. Для решения этой задачи используются модели
машинного обучения, такие как Cross Encoder~\cite{Nogueira_2019} и большие языковые модели (LLM).
Они позволяют реранжировать тексты на основе оценки наличия причинно-следственных, логических
и иных сложных связей.

Cross Encoder являются сравнительно небольшими моделями для попарной проверки текстов.
Эти модели принимают на вход два текста и совместно их обрабатывают с помощью cross-attention,
что позволяет выявить более тонкие логические связи и точнее провести фильтрацию.

Большие языковые модели
отлично понимают естественный язык и способны выявлять нетривиальные, многошаговые
логические цепочки, которые могли быть упущены при фильтрации через Cross Encoder.
Кроме того, LLM могут комбинировать фильтрацию с типизацией отношений,
что позволяет решить обе задачи за один запрос.
Однако такой подход является вычислительно сложным и, как следствие, дорогим, из-за чего
LLM используют только для классификации и дополнительной фильтрации уже отфильтрованных пар.

\subsection{Анализ существующих решений для построения графов знаний}

В настоящее время на рынке представлено множество решений
для управления знаниями, которые автоматизированное строят ГЗ.
Они варьируются от полноценных экосистем, предназначенных для корпоративного использования
до узконаправленных инструментов, решающих конкретные задачи.
Проведенный анализ позволил разделить решения на три группы, выявить ограничения и хорошие практики
существующих решений, а также сформировать базовый набор технологий,
применяемый при построении ГЗ в контексте систем управления знаниями.

\textbf{Готовые решения для корпоративного использования.}
В данную категорию включены полноценные экосистемы,
позволяющие автоматически трансформировать пользовательские текстовые данные
в графы знаний при помощью LLM.
Представителями этой категории являются такие решения как
Neo4j Graph Builder~\cite{neo4j_graph_builder} и Memgraph Unstructured2Graph~\cite{memgraph_u2g}.
Сами Neo4j и Memgraph являются графовыми базами данных,
предоставляющими широкий функционал «из коробки»,
в том числе возможность извлечения сущностей и связей из
пользовательских текстовых данных различных популярных форматов
для построения ГЗ на их основе.

Поскольку данная работа фокусируется на построении ГЗ,
в этих решениях будет изучен исключительно ответственный за это модуль.
Оба решения имеют открытый исходный код этих модулей.
За основу взят пайплайн Neo4j, однако логика Memgraph крайне схожа.
Изученный конвейер состоит из нескольких этапов:

\begin{enumerate}
      \item Парсинг и чанкинг. Исходные данные разбиваются на фрагменты размером
            4000 символов и перекрытием 200 символов [https://github.com/neo4j-labs/llm-graph-builder].
      \item Извлечение сущностей. Каждый чанк передается в LLM с инструкциями для извлечения
            троек \texttt{(Entity A — [relation type] — Entity B)}.
      \item Разрешение кореференций. Извлеченные имена сущностей
            дедуплицируются с помощью большой языковой модели
            для избежания конфликтов
            (например, сущности «ИИ» и «Искусственный интеллект» объединяются в одну с именем
            «Искусственный интеллект»).
      \item Запись данных. Обработанные данные записываются в базу данных.
            Они формируют два графа:
            лексический, сохраняющий последовательность фрагментов в файле, и семантический,
            содержащий связи между сущностями.
      \item Создание сообществ. Используются
            графовые алгоритмы (Louvain, PageRank) для поиска сообществ,
            что позволяет создавать кластерные узлы (хабы),
            увеличивающие эффективность навигации RAG.
\end{enumerate}

Neo4j и Memgraph используют LangChain~\cite{Chase_langchain} для работы с LLM,
что позволяет использовать множество провайдеров,
а также запущенные локально модели.

Данные системы обладают рядом преимуществ:
\begin{enumerate}
      \item Поддержка множества форматов текстовых данных,
            таких как PDF, HTML и DOCX.
      \item Построение двухуровневых графов: лексический,
            отражающий последовательность фрагментов текстов,
            и граф сущностей, соединенных типизированными связями.
            Это упрощает сбор всего содержимого одной заметки, например,
            в случае необходимости пересказа ее содержимого
      \item Формирование сообществ с описаниями,
            что экономит контекст при ответе на общие вопросы,
            затрагивающие множество заметок на общую тему.
\end{enumerate}

Однако они также имеют ряд критических недостатков, ставящих
под сомнение использование подобных решений для персональных баз данных.
\begin{enumerate}
      \item Отсутствие поддержки частичного обновления. Даже при минимальных модификациях
            файлов необходимо полностью перезапускать дорогостоящий пайплайн для построения ГЗ.
      \item Разрыв между данными и графом. Поскольку в этих сервисах отсутствуют
            механизмы синхронизации содержимого заметок, пользователю необходимо
            вручную загружать обновленные файлы на сервер базы данных,
            что увеличивает вероятность ошибок
            и делает процесс более трудоемким.
      \item Инфраструктурные издержки. Для обработки данных и взаимодействия с ГЗ необходима
            работающая база данных, что ведет к потреблению дополнительных вычислительных ресурсов,
            и, как следствие, дополнительным затратам, если сервис работает на выделенном сервере.
      \item Чрезмерное доверие к LLM. В пайплайнах этих решений отсутствует фильтрация кандидатов.
            В них вся работа по извлечению сущностей и анализу связей
            полностью лежит на больших языковых моделях, которым свойственны «галлюцинации».
            Это влечет создание бессмысленных и выдуманных связей,
            которые делают извлечение информации неэффективным.
\end{enumerate}

\textbf{Академические автоматизированные пайплайны.}
В качестве альтернативы решениям, опирающимся на
использование тяжеловесных LLM было найдено решение~\cite{Ain_2025},
эффективно использующее сравнительно небольшие модели и другие
методики для построения ГЗ.

Его пайплайн состоит их следующих этапов:
\begin{enumerate}
      \item Чанкинг. Файлы разбиваются на фрагменты, каждый из которых будет
            обрабатываться независимо, то есть для каждого чанка будет создан независимый ГЗ.
      \item Извлечение сущностей. Вместо больших языковый моделей данное решение использует
            компактную модель SqueezeBERT~\cite{Iandola_2020}. Процесс извлечения аналогичный:
            в модель передается фрагмент текста и инструкции.
      \item Привязка к онтологиям. Для стандартизации имен сущностей и снижения галлюцинаций
            все извлеченные сущности сверяются с базами данных, такими как DBpedia или Wikidata.
      \item Вычисление семантического сходства. Для создания связей между сущностями
            используется сравнение косинусного сходства их плотных эмбеддингов,
            созданных моделью Sentence-BERT~\cite{Reimers_2019}, с пороговым значением.
      \item Слияние графов. Все ГЗ, построенные для фрагментов
            текстов объединяются в единую сеть на основе общих узлов и связей.
\end{enumerate}

Ключевые достоинства этого подхода включают:
\begin{enumerate}
      \item Вычислительную эффективность.
            Использование небольших моделей машинного обучения сильно экономит
            вычислительные ресурсы и позволяет использовать решение полностью локально.
      \item Использование онтологий. Это решение упрощает этап разрешения кореференций,
            уменьшая вероятность галлюцинаций
            и убирая необходимость использования LLM для дедупликации.
\end{enumerate}

У данного решения есть несколько недостатков,
делающими этот подход неподходящим для решения главной проблемы этой работы.

\begin{enumerate}
      \item Обязательное участие человека. Поскольку авторам важно качество полученного ГЗ,
            сгенерированный граф тщательно корректируется человеком, что делает процесс
            лишь частично автоматизированным.
      \item Малое разнообразие сущностей. Использование жестких списков допустимых сущностей
            может повлечь ошибочную фильтрацию непопулярных терминов или аббревиатур.
\end{enumerate}

\textbf{Облачные приложения для заметок.}
Индустрия также имеет примеры приложений для ведения заметок,
которым удалось интегрировать ГЗ для увеличения эффективности их использования.

Облачные сервисы, такие как Mem.ai~\cite{Mem_ai}, являются
программным обеспечением как услугой (SaaS), полностью разгружая пользователя
от организации и связывания заметок.
Mem.ai, по аналогии с MemGraph и Neo4j, строит двухуровневый граф
(лексический слой,
сохраняющий последовательность фрагментов в файле, и семантический,
содержащий связи между сущностями).
Это обеспечивает эффективную навигацию алгоритмам RAG.
В данном сервисе ГЗ используются как для подсказок во время написания заметок,
так и для извлечения и синтеза информации посредством интегрированного чат-бота.

Несмотря на все удобства, у Mem.ai существуют недочеты,
касающиеся приватности и обработки данных:
\begin{enumerate}
      \item Закрытый исходный код. Пользователи не имеют доступа к логике построения ГЗ
            и обработки их заметок.
      \item Привязка к поставщику. Данные хранятся на серверах компании,
            что ставит под сомнение их безопасность,
            и делает их экспорт в случае закрытия сервиса затруднительным:
            сервис предоставляет возможность экспорта сырого текста заметок,
            но не ГЗ.
\end{enumerate}

\textbf{Плагины Obsidian.}
Приложение имеет поддержку пользовательских расширений,
которые добавляют новый функционал в программу.
По результатам исследования выяснилось,
что расширений, позволяющих автоматизировать процесс построения ГЗ не существует.
Несмотря на это, есть плагины-ассистенты, упрощающие некоторые его этапы.

\begin{enumerate}
      \item \textbf{Smart Connections}~\cite{smart_connections}. Плагин, использующих векторный поиск для
            формирования списка схожих заметок.
            Инструмент способен находить только семантически близкие заметки.
            Кроме того, плагин не проводит никакой дополнительной фильтрации результатов,
            а также не типизирует результаты, чем перекладывает
            большую часть работу на пользователя.
      \item \textbf{Graph Analysis}~\cite{graph_analysis}. Инструмент, подсказывающий связи на основе существующих.
            Данное решение опирается только на уже построенный граф для вычисления статистик,
            на основе которых предлагает пользователю новые связи или хабы.
            Данный плагин может хорошо работать
            как дополнение к автоматизированному пайплайну построения ГЗ,
            но не как его замена.
\end{enumerate}

\textbf{Выводы.}
Решения, представленные на рынке в данный момент, заставляют пользователей идти на компромиссы.
С одной стороны есть либо полноценные дорогостоящие и требовательные к ресурсам базы данных,
либо закрытые облачные решения для ведения заметок,
которые все же решают задачу автоматизированного построения ГЗ.
С другой стороны существуют открытые легковесные инструменты,
позволяющие упростить некоторые этапы построения,
но требующие вмешательства человека.

Решение с открытым исходным кодом, способное вычислительно эффективно находить,
фильтровать, типизировать и сохранять связи, а также поддерживающее работу полностью локально,
имело бы огромное преимущество на фоне изученных решений.
Плюсом будет возможность модификации построенного графа пользователем,
как и поддержка инкрементальных обновлений.

\subsection{Методология оценки качества}

Для объективной оценки разрабатываемого пайплайна и его сравнения с аналогами
необходимо создать методологию тестирования.
Процесс тестирования требует как подготовки валидационного датасета,
так и выбора подходящих метрик.

\subsubsection*{Реализация графа знаний в среде Obsidian и требования к данным}

Благодаря гибкости Obsidian и нативной поддержке ссылок,
создание ГЗ в нем будет выглядеть естественно.
В качестве сущностей будут выступать \texttt{Markdown}-файлы,
а в роли семантических связей~--- гиперссылки.
В данной работе будут также рассматриваться типизированные ссылки в Obsidian,
которые добавляет расширение Dataview.

Также необходимо установить ряд ограничений для сохранения каноничности графа,
а также упрощения работы с ним.
\begin{enumerate}
      \item Корпус текстов должен обладать высокой семантической связностью,
            то есть между файлами должны быть ярко выраженные семантические отношения
            (иерархические, причинно-следственные, функциональные и другие).
      \item Содержание каждого файла должно быть целостным и описательным.
      \item Атомарность сущностей: заметка должна представлять одну сущность.
      \item Пространство допустимых сущностей задается файлами в хранилище.
            При этом создание рефлексивных связей
            (когда заметка ссылается сама на себя) не допускается.
\end{enumerate}

Раскроем третье структурное ограничение подробнее.
Оно означает, что в содержании файла должна находиться информация только
об одной конкретной сущности, такой как личность, термин, событие, и так далее.
Содержание файла может включать в себя описание этой сущности, ее определение, историю,
факты, а также описание связей с другими сущностями.
Во-первых, такой подход делает хранилище максимально похожим на граф знаний.
Во-вторых, это убирает необходимость извлекать сущности,
так как файлы в хранилище уже формируют список допустимых сущностей.
В-третьих, данное ограничение соответствует концепции
атомарности заметок в Zettelkasten~\cite{Ahrens_2017},
где одна заметка содержит в себе ровно одну идею или концепцию.

\subsubsection*{Стратегия валидации}

Оценка такой комплексной системы от начала до конца
является некорректной, так как она может скрывать ошибки промежуточных этапов.
Например, если на этапе извлечения были получены нерелевантные чанки,
то этап типизации заведомо провалится, так как он не получит нужных данных.
В связи с этим валидация конвейера будет разделена на три этапа:
оценку извлечения контекста,
оценку типизации отношений,
а также оценку сгенерированного графа для использования в RAG-системах.
Такой подход даст возможность изолированного измерения качества
работы каждого модуля системы и его тонкой настройки,
что положительно скажется на итоговом качестве пайплайна.

\subsubsection*{Оценка извлечения контекста}

Оценку этапа извлечения (retrieval) можно сформулировать как задачу воспроизведения контекстов,
в которых упоминается сущность:
для каждой пары \((s, e)\), где \(s\)~--- файл-источник (\texttt{source\_file}),
а \(e\)~--- целевая сущность (\texttt{entity}),
система должна вернуть фрагмент из текста \(s\),
в достаточной мере совпадающий с эталонным контекстом из «золотого стандарта».
Извлеченный фрагмент должен полностью содержать эталонное предложение.

Важно заметить, что при больших чанках вероятность случайно захватить
нужное предложение выше.
Однако в реальных сценариях
извлекаемые фрагменты не могут быть слишком большими,
так как это приведет к размытию смыслов и потере деталей,
что негативно повлияет на качество извлечения.

Извлеченные фрагменты текста подразделяются на две категории:
\begin{itemize}
      \item Прямое (явное) упоминание.
            В данную группу входят все предложения,
            которые содержат полное или общепринятое краткое наименование.
            Например, предложение «Билл Гейтс соосновал Microsoft» явно содержит название компании.
      \item Кореферентное упоминание. Этот класс включает случаи,
            в которых прямое имя сущности заменено на
            другие языковые маркеры, требующие контекстного разрешения
            Например, в тексте «Microsoft была основана 4 апреля 1975.
            Билл Гейтс был одним из тех, кто ее соосновал» местоимение «ее» отсылает к Microsoft.
\end{itemize}

Для оценки необходимо создать специальный датасет, состоящий
из предложений с упоминаниями сущностей, списка упомянутых сущностей,
а также типа упоминания.

\subsubsection*{Оценка LLM-типизации}

Оценка типизации связей сводится к задаче классификации направленного
отношения для упорядоченной пары сущностей $(e_1, e_2)$.
При условии, что этап извлечения предоставил набор релевантных фрагментов,
большая языковая модель должна определить наличие связи от $e_1$ к $e_2$,
и присвоить ей один из представленных классов.
В случае присвоения недопустимого класса или
ответа модели о несвязанности сущностей
(классификации как отдельный класс \texttt{no\_relation}),
ссылка отбрасывается.

Для оценки качества классификации используются стандартные метрики~\cite{Manning_2008}:

\begin{itemize}
      \item Точность (Precision)~--- доля истинных связей среди
            всех предсказанных алгоритмом. Метрика важна для оценивания алгоритма на устойчивость к шуму.
            В случае создания множества ложноположительных связей граф засорится и потеряет свою ценность.
      \item Полнота (Recall)~--- доля истинных связей из датасета,
            которые алгоритм смог успешно идентифицировать и типизировать.
      \item Взвешенная F-мера (Weighted F\textsubscript{1}-score)~--- гармоническое
            среднее между точностью и полнотой. Эта метрика устойчива к дисбалансу в типах связей:
            зачастую в данных универсальные классы (такие как упоминания или дополнения)
            встречаются значительно чаще, чем, например, опровержение или противопоставление.
            Она будет использоваться в качестве основной оценки.
\end{itemize}

\subsubsection*{Полная оценка GraphRAG}

Для комплексной оценки ответов, сгенерированных LLM на основе контекста, предоставленного RAG,
используются специализированные фреймворки, такие как RAGAS~\cite{Es_2023}.
Их фундаментальной идеей является использование LLM как судей
для оценки ответов других языковых моделей,
что позволяет оценивать семантику ответов, а не использованные в них слова.

В данной работе будет использован фреймворк RAGAS.
Из метрик, предоставляемых им, будут подсчитываться:

\begin{itemize}
      \item Точность извлеченного контекста (Context Precision, CP).
            Данная метрика оценивает способность RAG систем
            извлекать релевантную информацию и ранжировать ее по мере соответствия.
            В RAGAS каждый извлеченный фрагмент текста оценивается
            на наличие и долю в нем нужной для ответа информации.
            Высокое значение CP говорит о предоставлении языковой модели
            релевантных данных с минимальной долей шума в них.
      \item Полнота извлеченного контекста (Context Recall, CR). Эта метрика оценивает RAG
            на умение собирать всю необходимую для ответа информацию.
            Для расчета этой метрики эталонный ответ
            разбивается на независимые атомарные утверждения, каждое из которых
            оценивается на возможность выведения его на основе извлеченных данных.
            Чем больше утверждений в эталонном ответе могут быть выведены,
            тем выше значение метрики.
      \item Достоверность ответа (Answer Faithfulness, F) направлена на измерение достоверности ответа,
            то есть степени,
            в которой сгенерированный ответ опирается исключительно на предоставленный контекст,
            а не на знания самой модели.
            Процесс оценки заключается в выделении всех тезисов в ответе и проверке
            возможности их подтверждения на основе предоставленных фрагментов текстов.
            Если модель использовала в ответе внешнюю информацию, то мера достоверности снижается.
      \item Корректность ответа (Answer Correctness, AC)
            измеряет семантическое и фактологическое соответствие
            сгенерированного ответа эталонному.
            Во фреймворке RAGAS она вычисляется как взвешенная сумма фактологического перекрытия
            (доля фактов, присутствующих как в ответе системы, так и в эталоне)
            и семантического сходства эмбеддингов двух ответов.
\end{itemize}

Использование такого набора метрик позволяет получить полную картину
работы пайплайна: от способности алгоритма GraphRAG предоставить релевантную информацию,
до способности модели использовать предоставленную информацию для генерации ответа.

\subsubsection*{Моделирование сложных сценариев для тестирования}

Чтобы проверить устойчивость алгоритма и его применимость в реальных условиях использования,
необходимо выявить наиболее сложные или важные ситуации,
с которыми встречаются алгоритмы построения ГЗ и GraphRAG.

\begin{itemize}
      \item Обнаружение сообществ. Сообщества~--- подграфы,
            вершины внутри которых плотно связаны между собой, но слабо с внешними.
            Они формируются из сущностей из одной тематики. Например,
            заметки с рецептами создают один кластер,
            а конспекты лекций по машинному обучению~--- другой.
            Современные алгоритмы ГЗ обычно автоматически
            их обнаруживают и создают для них отдельную вершину,
            в которой содержится краткая сводка~\cite{Edge_2024}.
            Эти вершины полезны при ответе на широкие вопросы,
            такие как
            «Какие темы из области математического анализа я изучал в этом месяце?»,
            а также для навигации алгоритмов поиска релевантных заметок.
            Они убирают нужду каждый запрос проходить по всем вершинам и
            обобщать их содержимое.
            Поскольку алгоритм, представленный в этой работе,
            не будет искать и создавать сообщества,
            уместно оценить возможность находить уже созданные.
            В рамках валидационного датасета вершины-сообщества будут созданы вручную.
      \item Проверка на устойчивость к шуму. В данном случае можно использовать омонимию,
            то есть слова или термины, совпадающие по написанию, но различные по смыслу.
            Такая структура заставит алгоритм использовать содержимое заметок для поиска и фильтрации.
            Также возможно создание графа из тесно связанных сущностей, что позволит
            протестировать способность алгоритма улавливать тонкие
            различия между типами связей в однородном тексте.
      \item Многошаговый поиск. Линейные цепочки из заметок, где
            каждая вытекает из предыдущей, проверяют
            эффективность поиска длинных логических путей.
\end{itemize}



\clearpage
\section{Разработка алгоритмов}
\subsection{Архитектура алгоритма построения графа знаний}

Разработанная система является конвейером для автоматического нахождения,
типизации и инкрементального сохранения семантических связей между
заметками (Markdown-файлами) в хранилище Obsidian.
Процесс построения состоит из нескольких шагов:
от сканирования хранилища до записи найденных отношений в файлы.
Алгоритм поддерживает сохранение состояния файлов в директории и
контрольные точки (checkpoints),
которые позволяют алгоритму продолжить работу с момента прерывания и
нацелены на экономию вычислительных ресурсов.
Алгоритм реализован на языке Python
(исходный код доступен в директории \texttt{src/kg\_builder}).
Общая схема конвейера представлена на рисунке~\ref{fig:pipeline}.

\begin{figure}[H]
    \centering
    \begin{tikzpicture}[
            node distance = 6mm and 30mm,
            block/.style = {
                    rectangle, rounded corners=3pt,
                    draw=black!60, fill=blue!8,
                    text width=72mm, minimum height=9mm,
                    align=center, font=\small
                },
            ioblock/.style = {
                    rectangle, rounded corners=3pt,
                    draw=black!40, fill=gray!10,
                    text width=72mm, minimum height=8mm,
                    align=center, font=\small\itshape
                },
            sideblock/.style = {
                    rectangle, rounded corners=3pt,
                    draw=black!50, fill=orange!12,
                    text width=36mm, minimum height=8mm,
                    align=center, font=\small
                },
            arrow/.style = {-Stealth, thick, black!70},
            dasharrow/.style = {Stealth-Stealth, dashed, thick, black!50},
            label/.style = {font=\scriptsize\itshape, text=black!60}
        ]

        % Main flow
        \node[ioblock]  (vault)    {Хранилище Markdown-файлов (vault)};
        \node[block,    below=of vault]  (scan)     {1. Обнаружение изменений\\{\scriptsize\normalfont VaultManager}};
        \node[block,    below=of scan]   (index)    {2. Разбивка на чанки и индексация\\{\scriptsize\normalfont VectorStore / LanceDB}};
        \node[block,    below=of index]  (retrieval){3. Гибридный поиск пар-кандидатов\\{\scriptsize\normalfont BM25 + векторный поиск}};
        \node[block,    below=of retrieval] (rerank){4. Переранжирование кандидатов\\{\scriptsize\normalfont Cross Encoder}};
        \node[block,    below=of rerank] (llm)      {5. Типизация связи\\{\scriptsize\normalfont LLM (локально или API)}};
        \node[block,    below=of llm]    (export)   {6. Запись связей в хранилище\\{\scriptsize\normalfont ExportService}};
        \node[ioblock,  below=of export] (result)   {Обновленные заметки с типизированными связями};

        % Side checkpoint block, aligned vertically between retrieval and llm
        \node[sideblock, right=of rerank] (meta)
        {Метаданные\\и чекпоинты\\{\scriptsize\normalfont MetadataManager}};

        % Main arrows
        \draw[arrow] (vault)     -- node[label, right]{markdown-файлы}    (scan);
        \draw[arrow] (scan)      -- node[label, right]{новые / обновленные файлы} (index);
        \draw[arrow] (index)     -- node[label, right]{чанки + эмбеддинги}  (retrieval);
        \draw[arrow] (retrieval) -- node[label, right]{пары-кандидаты}     (rerank);
        \draw[arrow] (rerank)    -- node[label, right]{отфильтрованные пары} (llm);
        \draw[arrow] (llm)       -- node[label, right]{тип связи + reasoning}(export);
        \draw[arrow] (export)    --                                           (result);

        % Dashed arrows to/from checkpoint block
        \draw[dasharrow] (meta.west) -- (rerank.east);
        \draw[dasharrow] ([yshift=4mm]meta.west) -- ([yshift=4mm]retrieval.east);
        \draw[dasharrow] ([yshift=-4mm]meta.west) -- ([yshift=-4mm]llm.east);

    \end{tikzpicture}
    \caption{Общая схема конвейера построения графа знаний}
    \label{fig:pipeline}
\end{figure}

\subsubsection*{Этап 1. Сканирование хранилища и обнаружение изменений (VaultManager)}
На этом этапе модуль \texttt{VaultManager} рекурсивно обходит директорию хранилища,
формируя список всех \texttt{markdown}-файлов.
Для каждого файла вычисляется SHA-256 хеш-сумма его содержимого для идентификации изменений.
Полученный хэш сравнивается с сохраненным с предыдущего запуска.
В результате файлы разбиваются на три категории: новые (отсутствующие в метаданных),
измененные (не совпадает хеш) и удаленные (есть в метаданных, но отсутствуют в директории).
Если метаданных нет, то все файлы считаются новыми.

Данный подход позволяет обрабатывать только измененные файлы,
а не все хранилище целиком.

\subsubsection*{Этап 2. Разбивка на чанки и индексация (VectorStore / LanceDB)}
Каждый новый или измененный файл проходит процесс индексации.
Эта процедура состоит из двух частей: обработки содержимого заметок и подготовки индексов.

Изначально текст каждого Markdown-файла может состоять из двух частей:
основного содержимого и блока ссылок, отделенного заголовком \texttt{link\_header},
если они были сгенерированы ранее.
Индексируется только основной текст заметки.

Текстовый индекс строится на основе лемматизированного содержимого файлов.
Лемматизация проводится при помощи модели \texttt{ru\_core\_news\_sm}
или \texttt{en\_core\_web\_sm} из \texttt{spaCy}~\cite{spacy_software}
для файлов с содержимым на русском или английском языке соответственно.
Из текстов извлекаются леммы всех алфавитных и числовых токенов.
Преобразованные тексты передаются в LanceDB~\cite{LanceDB_online}, где на них строится FTS-индекс,
использующий алгоритм BM25.
Данный подход позволяет увеличить точность поиска, нивелируя проблему, связанную с изменением форм слов,
в которых, например, могут быть изменены окончания или беглые гласные в корне,
что особенно распространено в русском языке.
Стоит отметить, что вместе с преобразованным текстом сохраняется и оригинальный.

Построение векторного индекса начинается с разбиения содержимого файла на перекрывающиеся
чанки. Используется рекурсивное посимвольное разбиение (\texttt{RecursiveCharacterTextSplitter} из LangChain).
Данная стратегия разбивает текст иерархически по заданным разделителям, расположенным в порядке их приоритета,
(символ новой строки, точка и т.д.), при этом не превышая заданный размер чанка
(по умолчанию 3,000 символов).
Эта стратегия сохраняет целостность предложений и структуру Markdown-файлов, что крайне важно
для вычисления эмбеддингов и формирования контекста для LLM.

Далее для каждого чанка следует вычисление плотного векторного
эмбеддинга при помощи модели из Sentence-Transformers~\cite{sentence_transformers_hf},
например популярной мультиязычной модели \texttt{BAAI/bge-m3}.

Полученные индексы объединяются друг с другом, а также с метаданными в виде пути к файлу и хеша,
после чего сохраняются в базу данных.
Хранение обоих форматов данных в одной таблице позволяет использовать гибридный поиск.
LanceDB позволяет регулировать баланс между векторным и полнотекстовым поиском при помощи параметра
\texttt{vector\_weight} (где $\texttt{vector\_weight} \in [0, 1]$): при \texttt{vector\_weight}~$= 1.0$ используется только ANN-поиск,
при \texttt{vector\_weight}~$= 0.0$~--- только BM25,
промежуточные значения используют оба поиска, сортируя результаты по унифицированной оценке релевантности.

\subsubsection*{Этап 3. Гибридный поиск пар-кандидатов (Retrieval)}

На данном этапе для каждого чанка выполняется поиск семантически близких чанков из других файлов.
Доступно три стратегии поиска:

\begin{itemize}
    \item \texttt{broad}, основанная на векторном поиске.
          Для каждого чанка текущего файла производится поиск $k$ ближайших фрагментов
          из других файлов.
    \item \texttt{strict}.
          Она основана на поиске прямых упоминаний сущностей в текстах заметок.
          Поиск производится по лемматизированному тексту.
    \item \texttt{combined} объединяет обе стратегии с дедупликацией. Результаты
          \texttt{broad} и \texttt{strict} проходят соответствующий им пайплайн извлечения и обработки,
          после чего фрагменты объединяются и дедуплицируются с помощью коэффициента Жаккара,
          вычисляемого по леммам.
          Результаты \texttt{strict} приоритизируются. То есть
          если фрагмент, извлеченный \texttt{strict}, будет иметь значение коэффициента с фрагментом,
          извлеченным \texttt{broad}, больше порога,
          то чанк от \texttt{broad} будет отброшен.
\end{itemize}


Особенности работы и недостатки каждой стратегии описаны в разделе 1.2.

Результатом поиска является \texttt{INITIAL\_RETRIEVAL\_K} фрагментов текста.

\subsubsection*{Этап 4. Переранжирование кандидатов (Cross-Encoder)}
Данный этап проходят только кандидаты, извлеченные векторным поиском,
поскольку среди них зачастую много ложноположительных.

Для снижения шума применяется дополнительная фильтрация кандидатов с помощью Cross Encoder.
Каждая пара (чанк источника, чанк цели) подается на вход Cross Encoder
(модель типа BAAI/bge-reranker-v2-m3~\cite{bge_reranker_v2})
с инструкцией-вопросом.
Cross Encoder вычисляет более точную оценку семантической связности.
Пары сортируются по убыванию оценки. В финальный список попадают \texttt{RERANKER\_TOP\_K} пар,
у каждой из которых оценка не ниже порогового значения.

Все прошедшие фильтрацию чанки одной заметки группируются по файлу-цели.
На их основе создаются структуры \texttt{CandidatePair},
содержащие исходный и целевые чанки, а также пути к соответствующим файлам и оценку близости
из LanceDB и Cross Encoder.

Результатом этапа является список структур \texttt{CandidatePair},
к которым добавлены результаты лексического поиска, если он был использован.

\subsubsection*{Этап 5. Типизация связи (LLM)}
Прошедшие фильтрацию сгруппированные пары-кандидаты
передаются в языковую модель для определения точного типа семантической связи.
Для каждой такой пары отбираются не более \texttt{LLM\_TOP\_K\_CTX}
наиболее релевантных чанков для формирования контекста
на основе оценки Cross Encoder.
Важно уточнить, что результаты стратегии \texttt{strict}
являются более приоритетными, а
результаты стратегии \texttt{broad}
служат для дополнения выборки и
добавляются в итоговый список
по остаточному принципу на основе их оценок.

Логика запросов к LLM поддерживает параллелизм и батчинг.
Количество одновременных запросов к модели ограничено заданным числом и контролируется семафором.
Кроме того, в рамках одного запроса может проводиться типизация сразу нескольких ссылок.
Это значительно снижает накладные расходы и время обработки.

Важно отметить обработку ошибок: у каждого батча есть три попытки, при израсходовании которых
батч откладывается и логируется ошибка. Он будет повторно обработан при следующем запуске.

В рамках запроса модель получает системный промпт со списком допустимых типов связей и инструкциями,
а также пользовательский промпт с парами файлов и их контекстом.

Ответ модели возвращается в формате JSON. Его парсинг производится посредством JsonOutputParser
из LangChain, позволяющий восстанавливать синтаксически некорректный JSON.
При возврате моделью раздела с размышлениями
он логируется и вырезается для упрощения парсинга.

Для работы с LLM используется LangChain,
который позволяет интегрировать
как локальные модели
(например, запущенные через llama.cpp~\cite{Gerganov_llamacpp}),
так и облачные провайдеры (OpenAI, Google, и другие).

После каждого батча предсказания модели сохраняются в файл,
формируя контрольную точку, что позволяет возобновить типизацию
в случае ошибки или аварийного завершения программы.
Данная оптимизация призвана
убрать необходимость повторять выполненную работу.

\subsubsection*{Этап 6. Сохранение связей (ExportService)}
На заключительном этапе типизированные связи сохраняются в одном из трех режимов:

\begin{itemize}
    \item \texttt{inplace}.
          Связи добавляются непосредственно в исходные Markdown-файлы под специальным заголовком
          \texttt{LINK\_HEADER}.
          При наличии связей с предыдущих запусков
          новые связи добавятся к уже имеющимся, не создавая дубликатов.
    \item \texttt{json}. Результаты экспортируются в JSON-файл без изменения исходных заметок.
          Данный режим создан для удобного переиспользования результатов, например, для вычисления метрик.
    \item \texttt{export}. Создается полная копия хранилища с вписанными связями,
          исходные файлы остаются нетронутыми.
          Этот формат сохранения полезен для экспериментов,
          таких как оценка GraphRAG, требующих полноценного хранилища для работы.
\end{itemize}

Формат ссылки конфигурируется
(по умолчанию используется синтаксис Dataview).

\subsubsection*{Управление состоянием и отказоустойчивость (MetadataManager)}

При первом запуске конвейера стандартные параметры из \texttt{config.py}
сохраняются в \texttt{config.json} внутри директории \texttt{.kg\_builder/}.
Последующие запуски считывают этот файл и выставляют настройки, заданные в нем,
что обеспечивает воспроизводимость и
повышает удобство использования на разных хранилищах с индивидуальными параметрами.

Опция \texttt{--ignore-local-config} заставляет пайплайн
проигнорировать \texttt{config.json} и использовать значения из \texttt{config.py}.

Алгоритм сохраняет свое состояние с помощью создания снимка, который включает
текущий этап,
его прогресс и прочие метаданные, необходимые для возобновления работы.

Наибольшей детализацией сохранения обладают этапы ранжирования кандидатов и LLM-классификации отношений.
У них есть логика контрольных точек, при которых сохраняется список уже обработанных файлов или пар-кандидатов.
При возобновлении работы пайплайна на одном из этих шагов в алгоритм будут переданы только необработанные данные.

При последующих запусках \texttt{KnowledgeGraphBuilder} проверяет последнее сохраненное состояние.
Если предыдущий запуск был прерван, то он продолжит обработку с того же шага, минуя уже выполненные.

Это сохраняет значительное количество вычислительных ресурсов и времени, особенно на дорогостоящем LLM-этапе.

Все метаданные сохраняются в отдельной директории \texttt{.kg\_builder/.out/}:
\begin{itemize}
    \item \texttt{candidates.json}~--- начальные кандидаты до переранжирования;
    \item \texttt{reranked\_candidates.json}~--- кандидаты после переранжирования,
          обогащенные результатами LLM (поля \texttt{llm\_type}, \texttt{llm\_reasoning}
          дописываются после классификации)
    \item \texttt{predictions\_partial.json}~--- предсказания LLM, обновляемые инкрементально
    \item \texttt{run\_state.json}~--- текущее состояние алгоритма
\end{itemize}

Специальный флаг \texttt{--fresh-start} позволяет перезаписать метаданные,
в том числе конфиг и список обработанных файлов.

\subsection{Алгоритм GraphRAG}

В качестве референсного решения был выбран ObsidianRAG~\cite{obsidianrag}.
Его алгоритм состоит из нескольких шагов:
\begin{enumerate}
      \item Индексация: содержимое \texttt{markdown}-файлов разбивается на чанки
            разметом 1500 символом с перекрытием в 300 символов.
            Полученные фрагменты вместе с извлеченными из текста ссылками
            записываются в ChromaDB - векторную базу данных -
            которая создает для них плотные векторные эмбеддинги.
      \item Первоначальный поиск: запрос пользователя обрабатывается
            с помощью \texttt{EnsembleRetriever}, использующим BM25 и ANN
            с весами 0.4 и 0.6 соответственно.
            Извлекается N наиболее релевантных заметок.
      \item Переранжирование: полученные файлы передаются в Cross Encoder,
            например, \texttt{BAAI/bge-reranker-v2-m3}.
            В результате фильтрации остается 6 наиболее близких фрагментов.
      \item Расширение контекста: результат дополняется пятью файлами,
            связанными ссылками с оставшимися.
      \item Генерация ответа: собранное содержимое заметок
            передается в локальную LLM (Ollama).
\end{enumerate}

У данного пайплайна есть ряд недостатков:

\begin{enumerate}
      \item Одноэтапное расширение контекста соседями.
            При таком подходе теряются сложные связи, требующие многошаговых рассуждений.
            А также возрастает вероятность остаться в семантической ловушке: ситуации,
            когда новая информация не может быть извлечена
            из-за сильной разницы в эмбеддингах и ключевых словах.
            Суть ссылок заключается в решении этой проблемы, но автор ими пренебрегает.
      \item Сбор контекста без учета его релевантности.
            В проекте уже используется Cross Encoder для первоначального поиска,
            его можно легко переиспользовать для фильтрации фрагментов, не относящихся к запросу.
\end{enumerate}

На основе этого решения был разработан улучшенный четырехэтапный пайплайн,
поддерживающий типизированные ссылки и многошаговое расширение контекста.
Он нацелен на решение фундаментальных проблем эталонного решения.
Алгоритм реализован на языке Python
(исходный код доступен в папке \texttt{src/graphrag}).

\begin{figure}[H]
      \centering
      \begin{tikzpicture}[
                  node distance = 6mm and 32mm,
                  block/.style = {
                              rectangle, rounded corners=3pt,
                              draw=black!60, fill=blue!8,
                              text width=76mm, minimum height=9mm,
                              align=center, font=\small
                        },
                  ioblock/.style = {
                              rectangle, rounded corners=3pt,
                              draw=black!40, fill=gray!10,
                              text width=76mm, minimum height=8mm,
                              align=center, font=\small\itshape
                        },
                  sideblock/.style = {
                              rectangle, rounded corners=3pt,
                              draw=black!50, fill=orange!12,
                              text width=38mm, minimum height=8mm,
                              align=center, font=\small
                        },
                  arrow/.style  = {-Stealth, thick, black!70},
                  dasharrow/.style = {dashed, thick, black!50, -Stealth},
                  label/.style  = {font=\scriptsize\itshape, text=black!60}
            ]

            % Main flow
            \node[ioblock]  (query)   {Запрос пользователя};
            \node[block, below=of query]   (seed)
            {1. Первоначальный поиск\\
            {\scriptsize\normalfont Гибридный поиск (BM25~+~вектор)}};
            \node[block, below=of seed]    (trav)
            {2. Обход соседей\\
            {\scriptsize\normalfont базовый или типизированный BFS}};
            \node[block, below=of trav]    (prune)
            {3. Отбор контекста\\
            {\scriptsize\normalfont Чанкинг, переранжирование}};
            \node[block, below=of prune]   (gen)
            {4. Генерация ответа\\
            {\scriptsize\normalfont LLM, структурированный промпт}};
            \node[ioblock, below=of gen]   (answer)  {Ответ};

            % Side block
            \node[sideblock, right=of trav] (kg)
            {Граф знаний\\{\scriptsize \texttt{rel\_type:: [[T]]}}};

            % Main arrows
            \draw[arrow] (query)  -- node[label, right]{} (seed);
            \draw[arrow] (seed)   -- node[label, right]{TOP\_K\_SEED заметок} (trav);
            \draw[arrow] (trav)   -- node[label, right]{заметки, дополненные соседями} (prune);
            \draw[arrow] (prune)  -- node[label, right]{TOP\_K\_CONTEXT релевантных чанков} (gen);
            \draw[arrow] (gen)    -- (answer);

            % Dashed arrow from KG to traversal
            \draw[dasharrow] (kg.west) -- node[label, above]{\scriptsize исходящие ссылки} (trav.east);

      \end{tikzpicture}
      \caption{Общая схема конвейера GraphRAG}
      \label{fig:graphrag_pipeline}
\end{figure}

\textbf{Этап 1. Первоначальный поиск.}
На этом этапе выполняется гибридный поиск по LanceDB.
Он включает в себя ANN-поиск по плотным эмбеддингам и BM25-поиск по
лемматизированному тексту. На его основе собирается $TOP\_K\_SEED$ (по умолчанию~$= 5$)
наиболее близких заметок-кандидатов. Они добавляются в множество посещенных файлов.

В отличие от референсного решения,
на данном этапе не выполняется переранжирование заметок с помощью Cross Encoder.
Оно было перенесено на шаг 3.
Это решение было принято исходя из точности алгоритма извлечения, а также для
смягчения критериев для отбора заметок,
что потенциально может увеличить полноту собранного контекста.

\textbf{Этап 2. Обход соседей.}

Расширение контекста может быть проведено двумя способами:

\begin{itemize}
      \item Базовый обход в ширину. Для каждой заметки формируется список исходящих ссылок.
            Типы ссылок при этом игнорируются, как и релевантность их содержимого.
            Все непосещенные файлы, соединенные с текущим, добавляются в очередь для обхода.
      \item Типизированный обход в ширину. Для каждой заметки, связанной с текущей,
            формируется запрос в Cross Encoder модель формата:
            \begin{center}
                  \texttt{relation: target. summary}
            \end{center}
            где \texttt{relation}~--- тип связи, \texttt{target}~--- имя целевой заметки,
            \texttt{summary}~--- ее краткое содержание (первый абзац).
            Cross Encoder вычисляет оценку близости соседней заметки с текущей.
            Все кандидаты с оценкой ниже пороговой не добавляются в очередь обхода.
            Это позволяет фильтровать заметки на этапе сбора контекста, что
            может снизить количество шума, но также уменьшить полноту контекста из-за раннего отброса.
\end{itemize}

Результатом этапа является множество потенциально релевантных заметок.

Наглядное сравнение стратегий обхода представлено в таблице~\ref{tab:traversal_compare}.

\begin{table}[H]
      \centering
      \begin{tabularx}{\textwidth}{|l|X|X|}
            \hline
            \textbf{Критерий}             & \textbf{Base}    & \textbf{Typed}                     \\
            \hline
            Тип связи                     & Игнорируется     & Включается в оценку                \\
            \hline
            Оценка кандидатов             & Отсутствует      & Cross Encoder                      \\
            \hline
            Текст для оценки              & ---              & \texttt{relation: target. summary} \\
            \hline
            Фильтрация                    & Нет (все соседи) & Пороговая по оценке                \\
            \hline
            Ограничение числа новых узлов & Нет              & Нет                                \\
            \hline
      \end{tabularx}
      \caption{Сравнение стратегий обхода графа}
      \label{tab:traversal_compare}
\end{table}

\textbf{Этап 3. Отбор содержимого.}
Содержимое каждой заметки, отобранной на предыдущем этапе, разбивается на фрагменты
с помощью того же сплиттера, что используется при индексации хранилища (глава 2.1).

Каждый полученный чанк получается оценку Cross Encoder схожести с запросом пользователя.
Сохраняется $TOP\_K\_CONTEXT$ (по умолчанию 10) наиболее релевантных фрагментов,
которые сформируют итоговый контекст и будут отправлены в LLM.

В отличие от референсного решения, в котором фильтрация происходит на этапе сбора контекста,
здесь происходит отбор всего собранного содержимого заметок после обхода графа.
Это позволяет увеличить полноту выборки,
но допускает наличие в ней большего количества нерелевантных фрагментов.

\textbf{Этап 4. Генерация ответа.}
Запрос в LLM формируется на основе промпта, который дает модели необходимые инструкции
(например, отвечать на языке пользователя, несмотря на язык предоставленного контекста,
и опираться только на предоставленные данные),
сообщения пользователя и
отобранных фрагментов текстов \texttt{markdown}-файлов,
разграниченные специальным разделителем.

Предсказание модели преобразуется в  формат JSON с помощью
\texttt{JsonOutputParser} из LangChain.
В ответе ожидается два поля:
\texttt{reasoning}~--- рассуждения модели, и
\texttt{answer}~--- финальный ответ, возвращаемый пользователю.
В случае неудачного парсинга возвращается сырой ответ модели.


\section{Сбор и подготовка данных}

Для оценки качества работы пайплайна и возможности
сравнения его с аналогами необходимо создать репрезентативный датасет,
который бы удовлетворял требованиям, описанным в главе 1.4.

\subsection{Поиск данных}

В качестве первичной гипотезы рассматривалось использование существующих открытых
пользовательских хранилищ в формате Obsidian с их последующей модификацией и фильтрацией.
В ходе поиска было отобрано несколько качественных репозиториев:
\begin{itemize}
      \item личные заметки общего назначения~\cite{noodleslove_notes};
      \item структурированная документация коллектива \textit{Hundred Rabbits}~\cite{100r_co};
      \item корпус кулинарных рецептов \textit{based.cooking}~\cite{based_cooking};
      \item заметки по программной инженерии~\cite{keyvanakbary_notes}.
\end{itemize}

После их детального анализа выявился один общий критический недостаток:
нарушалось правило ``заметка = сущность''.
Наиболее распространенной причиной его нарушения были
склейки текстов на разные тематики в одном файле.

Помимо структурного несоответствия был выявлен ряд дополнительных недостатков:
\begin{itemize}
      \item Семантическая разреженность данных.
            Некоторые хранилища состояли в основном из дневниковых записей, заметок или рецептов,
            между которыми практически не было связей, либо они были неоднозначными.
      \item Несоответствие объема.
            Почти все найденные хранилища были либо слишком малы, либо слишком велики.
            К репозиториям с недостающим объемом пришлось бы искусственно добавлять данные, а из
            слишком больших удалять,
            что потенциально может создать массу висячих ссылок и
            спровоцировать потерю ценных связей и контекста.
      \item В некоторых репозиториях отсутствует явная указанная лицензия
            либо присутствуют фрагменты текстов или кода,
            потенциально нарушающие авторские права.
            Это делает обработку этих данных рискованной.
\end{itemize}

\textbf{Выбор Википедии как основного источника данных.}
В качестве следующего источника данных было решено рассмотреть Википедию.
У нее есть ряд важных преимуществ:
\begin{itemize}
      \item Естественная атомарность статей. Зачастую статьи соответствуют одной сущности,
            что согласуется с правилом ``одна заметка = одна сущность''.
      \item Наличие ссылок. В статьях изначально есть гиперссылки на другие статьи.
            Эти ссылки образуют граф,
            к которому можно получить доступ и использовать его как для сбора данных,
            так и для первоначальной разметки.
      \item Качество содержимого. Подавляющая часть популярных статей
            написана в энциклопедическом стиле,
            обладает высокой плотностью информации, ее последовательным и связным изложением.
\end{itemize}

Все вышеперечисленное делает Википедию
подходящим источником данных для валидационного датасета.

\subsection{Создание валидационного датасета}

\textbf{Построение датасета и восстановление связей.}
Формирование базового датасета проводилось при помощи Wikipedia API,
обращение к которому производилось посредством библиотеки \texttt{wikipediaapi}.
Данный инструмент позволяет сохранять тексты статей, а также предоставляет доступ к графу из ссылок.
Важно отметить, что граф поддерживает типизацию ребер, что также поможет формированию датасета.
Однако лишь малая его часть размечена, и в основном это ссылки, соединяющие популярные статьи.

Важно отдельно отметить, что статьи имеют специальный раздел, содержащий краткую сводку (summary).
Он будет использован качестве контекста пайплайном GraphRAG
для фильтрации заметок на этапе сбора контекста.

\textbf{Формирование структуры элементов датасета.}
Поскольку ресурсы ограничены,
необходимо определить несколько наиболее показательных сценариев для оценки алгоритма.
На основе анализа особенностей работы GraphRAG, описанных в главе 1.2,
был составлен список наиболее показательных структур
элементов валидационного датасета
и их содержимого:

\begin{itemize}
      \item \textbf{Hub-n-spoke (Hub)}. Структура в форме граф в виде звезды,
            состоящая из статьи-хаба и ее соседей.
            Главная задача алгоритма в этом случае - идентификация хаба.
            Алгоритмы GraphRAG зачастую пытаются обнаружить именно такие узлы,
            так как они дают непосредственный доступ к большому количеству информации.
            Тематика датасета -- производство полупроводников и микроэлектроника.
            Итоговый набор статей включает как общие концепции (закон Мура, фотолитография, кремний),
            так и конкретные технологии или компании (ASML, TSMC, транзисторы, архитектура ARM, Apple M1).
      \item \textbf{Versus}. Структура, состоящая из двух разных тематик, которые местами могут показаться схожими.
            Например, Apple как компания и apple как яблоко.
            Такой подход заставит алгоритм работать с содержимым заметок, а не только с названиями сущностей.
            В собранном наборе данных эта структура реализована на примере явления омонимии:
            корпорация Apple (включая связанные узлы, например, Стив Джобс, iPhone, macOS)
            противопоставляется яблоку как фрукту (яблочный пирог, сорт Грэнни Смит)
            и связанным с ним историческим фактам (Исаак Ньютон).
      \item \textbf{Realistic}. Данный элемент будет содержать две схожие тематики.
            Его задача - проверка созданных алгоритмом связей на шум.
            Для датасета было решено связать
            концепции искусственного интеллекта (перцептроны, машинное обучение,
            обратное распространение ошибки и т.д.) с их историей.
      \item \textbf{Sequence}. Образец,
            состоящий из последовательных или вытекающих друг из друга событий или терминов.
            В качестве предметной области для этого набора была выбрана хронология развития представлений
            о строении атома (от открытия катодных лучей и электрона до пудинговой модели,
            экспериментов Резерфорда и модели Бора).
            Этот подкорпус будет нацелен на проверку способности алгоритма находить сложные, многошаговые связи.
      \item \textbf{Personal}. Элемент,
            созданный в ручную по аналогии с ежедневными заметками и содержащий различные события,
            планы, наброски идей.
            В нем смоделированы типичные личные записи: привязки к датам,
            повседневные задачи (починка велосипеда, покупки), встречи (с научным руководителем, знакомыми)
            и проекты (работа над дипломом, исследование малых LLM).
            Данный образец поможет оценить работу в условиях малого и специфичного контекста,
            отличного от текстов статей.
\end{itemize}

Визуализации ГЗ,
созданных на основе этих подкорпусов,
изображена на рисунке~\ref{fig:all_dataset_graphs}.

\subsubsection{Алгоритм генерации связей}

Для формирования датасета был написан скрипт для выгрузки и
обработки статей с использованием \texttt{wikipediaapi},
поддерживающий два режима сохранения:

\begin{itemize}
      \item выгрузка заданного набора статей
      \item обход соседей заданной статьи в ширину с контролем глубины обхода
\end{itemize}

Во второй стратегии сохранения
соседи статьи ранжируются с помощью Cross Encoder модели на основе их краткой сводки
и отбирается \textit{N} наиболее релевантных,
что предотвращает уход алгоритма от тематики корневой статьи.
Далее алгоритм увеличивает глубину рекурсии и аналогично обрабатывает каждого отобранного соседа
до тех пор пока либо не сохранит заданное число статей либо не обойдет все статьи до заданной глубины рекурсии.

Каждая статья, прошедшая проверку Cross Encoder подвергается обработке.
В нее входит:
\begin{itemize}
      \item Нормализация названий файлов, включающая замену пробелов на подчеркивания и
            удаление некоторых служебных символов для
            однозначного соответствия названий статей названиям файлов.
      \item Очистка текста от сносок и маркеров (например, \texttt{[1]}),
            а также удаление служебных секций статьи, которые не несут
            содержательной информации для GraphRAG
            (например, ``См. также'', ``Литература'', ``Ссылки'').
      \item Формирование специального блока ссылок, извлеченных из графа Википедии, на другие статьи.
            В данном блоке остаются ссылки только на сохраненные статьи,
            что решает проблему висячих ссылок.
\end{itemize}

После сохранения статей начинается этап их ручной доразметки.
Этот шаг включает:
\begin{itemize}
      \item Анализ и фильтрация ссылок, сохраненных из графа Википедии.
            Все ссылки, для идентификации которых нет контекста в содержимом статьи, удаляются.
            Например, в статье Билл Гейтс сказано, что он соосновал Microsoft, следовательно связь
            Билл Гейтс -> (соосновал) Microsoft должна быть оставлена.
      \item Анализ типов ссылок. Как было упомянуто ранее,
            в графе Википедии часть ссылок уже типизирована.
            Эти типы дополняются недостающими, а также абстрагируются для сокращения их общего числа.
            Критично обеспечить оптимальный уровень абстракции,
            чтобы не запутать языковая модель:
            большое количество узких по смыслу ссылок может показаться лучшим решением,
            но на практике модель начнет выбирать неверные типы, например из-за недостатка контекста.
\end{itemize}

Для упрощения разметки и алгоритмов
все типы отношений
фиксируются на английском языке
(например, \texttt{Uses}, \texttt{Part of}, \texttt{Mentions}),
несмотря на то, что содержимое
некоторых заметок на русском языке.

В каждом подкорпусе присутствует файл .meta,
в котором содержится список статей с Википедии,
на основе которых он был создан,
а также список допустимых типов отношений.

Результатом ручной доразметки является блок классифицированных ссылок
в каждой статье,
оформленный в синтаксисе плагина Dataview:

\begin{verbatim}
## Related Connections:
- Mentions:: [[ASML_Holding]]
- Manufactures:: [[Integrated_circuit]]
- Sustains:: [[Moores_law]]
- Is a:: [[Photolithography]]
\end{verbatim}

\begin{table}[H]
      \centering
      \begin{tabular}{lcccc}
            \hline
            Подкорпус              & \# заметок & \# ссылок & \# типов связей & Ср. исх. ссылок на заметку \\
            \hline
            \texttt{hub\_n\_spoke} & 15         & 77        & 12              & 5.13                       \\
            \texttt{realistic}     & 15         & 79        & 12              & 5.27                       \\
            \texttt{versus}        & 15         & 44        & 14              & 2.93                       \\
            \texttt{sequence}      & 7          & 23        & 7               & 3.29                       \\
            \texttt{personal}      & 11         & 23        & 7               & 2.09                       \\
            \hline
            \textbf{Итого}         & 63         & 246       & ---             & 3.90                       \\
            \hline
      \end{tabular}
      \caption{Сводная статистика по эталонных подкорпусах}
      \label{tab:dataset_samples}
\end{table}


В таблице~\ref{tab:dataset_samples} приведена информация о собранных данных.
\texttt{hub\_n\_spoke} и \texttt{realistic} обладают самым высоким средним числом связей на заметку
(в среднем 5.2 исходящих ссылки),
то есть в них много семантически близких и связных сущностей.
Они помогают оценить пайплайн построения графа на устойчивость к шуму.
\texttt{versus}, несмотря на идентичное число заметок, имеет вдвое меньше
связей, поскольку в нем часть сущностей не имеют никаких семантических отношений.
Аналогичная ситуация в \texttt{sequence}: в нем изложена последовательность,
в которой каждая сущность в основном связана только с предыдущей и следующей.

Датасеты хранятся в директории \texttt{data/test\_vaults/}.


\subsubsection{Алгоритм извлечения контекста}

Для оценки retrieval-этапа был создан отдельный датасет,
состоящий из фрагментов текстов, в которых явно или косвенно
упоминаются сущности и подкорпуса.

В качестве минимального фрагмента текста было решено взять предложение,
содержащие упоминание. Это позволит сохранить контекст,
а также проверять алгоритм на точность извлечения,
то есть долю других предложений, не содержащих упоминание.

Для каждого элемента датасета было проведено
ручное извлечение всех фрагментов с упоминаниями.
На их основе создан набор JSON-файлов
(по одному на элемент эталонного набора данных для оценки алгоритма генерации ссылок),
хранящийся в директории \texttt{data/test\_retrieval/gold/}.
Каждый файл содержит массив кортежей \((s, e, sent, is\_direct)\),
где \(s\) -- файл-источник,
\(e\) -- целевая сущность,
\(sent\) -- предложение, в котором упомянута сущность,
\(is\_direct\) -- содержит ли предложение прямое упоминание сущности.
Допускается, что для одной пары
\((s,e)\) может присутствовать несколько
записей с разными извлеченными предложениями.

Важно уточнить, что для соблюдения требований к данным,
описанным в главе 1.4,
в датасете используются только чанки,
содержащие упоминания сущностей, представленных в подкорпусе.

\subsubsection{Алгоритм LLM-классификации отношений}

Для валидации LLM-классификатора отношений используется дополненная
версия описанного выше датасета для оценки retrieval.
В каждую запись в JSON-файлах было добавлено поле
\texttt{link\_type}, соответствующее истинному типу отношения.
Таким образом финальная структура кортежей имеет следующий вид:
\((s, e, sent, is\_direct, link\_type)\).

Значения \(link\_type\) были взяты из блока ссылок
датасета для оценки алгоритма генерации связей.
В случае, если сущность упомянута в нескольких
предложениях, каждое из них получает одинаковое значение
\(link\_type\).

\textbf{Вывод.}

Для оценки качества и устойчивости алгоритмов к разным сценариям использования были сформированы
датасеты, различающиеся тематикой, плотностью информацией, сложностью и набором сущностей,
для оценки как всего пайплайна, так и для конкретных его этапов.
Данный подход задуман для упрощения тонкой настройки гиперпараметров и
выбора оптимальных стратегий извлечения контекста и классификации связей.

Выбор в пользу Википедии обеспечил выполнение правила ``одна заметка = одна сущность''
и позволил получить управляемый объем данных с качественным содержимым, достаточным для оценки GraphRAG.



% \textbf{Требования к данным.}
% С учетом устройства предлагаемого конвейера (см. раздел~2.1) к датасету были предъявлены следующие требования:
% \begin{enumerate}
%     \item \textit{Структурная согласованность.} Основная единица хранения и обработки в системе~--- Markdown-файл (``заметка'').
%           На уровне графа знаний вершина однозначно идентифицируется путем файла,
%           поэтому требуется соблюдение принципа ``одна заметка = одна сущность'' (атомарность представления).
%     \item \textit{Наличие смысловых связей.} Корпус должен содержать нетривиальные межтекстовые отношения
%           (иерархические, причинно-следственные, функциональные и др.), иначе оценка качества типизации
%           становится статистически неустойчивой и малоинформативной.
%     \item \textit{Контролируемый объем.} Для построения ``золотого стандарта''
%           необходима ручная или экспертно контролируемая разметка.
%           Следовательно, требуется ограниченный по размеру подграф,
%           в котором разметка выполнима в рамках дипломного исследования.
%     \item \textit{Качество и связность текста.} Языковые модели и модели ранжирования устойчивее работают на связных описательных фрагментах.
%           Преобладание обрывочных записей, фрагментов кода и цитат без контекста приводит к росту доли ложноположительных и ложноотрицательных связей.
%     \item \textit{Юридическая воспроизводимость.} Источник должен иметь ясный лицензионный статус,
%           допускающий обработку и публикацию производных данных (в частности, размеченных подмножеств и статистик).
% \end{enumerate}

% \textbf{Ключевое ограничение ``заметка = сущность''.}
% В предлагаемой архитектуре идентификатором сущности выступает файл (его относительный путь в хранилище).
% Это допущение используется системно: при индексировании и разбиении на чанки каждый фрагмент текста наследует идентификатор файла;
% при генерации кандидатных связей сопоставляются пары ``файл-источник~--- файл-цель'';
% при экспорте результаты записываются обратно в Markdown в виде типизированных ссылок вида \texttt{Relation:: [[TargetNote]]}.
% Если одна заметка содержит множество разнотипных сущностей (например, дневниковую запись с перечислением задач, людей и событий)
% или, наоборот, одна сущность ``распылена'' по нескольким файлам, то возникает неоднозначность:
% \begin{itemize}
%     \item невозможно однозначно интерпретировать ребро графа как отношение между сущностями (ребро связывает \textit{документы}, а не \textit{концепты});
%     \item при разметке и оценке качества предсказаний существенно возрастает субъективность: эксперт вынужден решать, к какой именно сущности внутри текста относится классифицируемая связь;
%     \item рефакторинг исходного корпуса (дробление лонгридов, переименование, перенос фрагментов, нормализация ссылок) фактически превращается в создание нового датасета и искажает ``естественную'' топологию связей, которую требуется моделировать.
% \end{itemize}
% Таким образом, нарушение принципа атомарности в исходном хранилище является не частным неудобством, а фундаментальным методологическим препятствием для построения валидируемого ``золотого стандарта''.

% \textbf{Предварительный анализ открытых Obsidian-хранилищ.}
% В качестве первичной гипотезы рассматривалось использование существующих пользовательских хранилищ (vaults) формата Obsidian, опубликованных в open-source. На этапе отбора наличие или отсутствие уже проставленных ссылок не считалось критичным: предполагалось, что связи можно извлечь и типизировать путем ручной или полуавтоматической разметки.

% В ходе поиска были отобраны и проанализированы несколько публичных репозиториев:
% \begin{itemize}
%     \item \url{https://github.com/noodleslove/notes/} (личные заметки общего назначения);
%     \item \url{https://github.com/hundredrabbits/100r.co} (структурированная документация коллектива);
%     \item \url{https://github.com/lukesmithxyz/based.cooking} (корпус кулинарных рецептов);
%     \item \url{https://github.com/keyvanakbary/learning-notes} (заметки по программной инженерии).
% \end{itemize}

% Детальный анализ показал, что такие хранилища систематически не удовлетворяют ключевому требованию атомарности ``заметка = сущность''. 
% Типичные паттерны нарушений включают: дневниковые записи с множеством сущностей в одном документе, ``оглавления'' и MOC-страницы (map-of-content), 
% лонгриды с несколькими связанными темами, а также файлы, преимущественно состоящие из фрагментов кода и выписок. 
% Адаптация подобных хранилищ потребовала бы радикальной переработки: разбиения документов на атомарные сущности, переименования, 
% восстановления и перенормализации ссылок, а иногда полного переписывания содержания, что противоречит цели исследования 
% (оценивать алгоритм на реалистичных данных, а не на искусственно подготовленном корпусе).

% Помимо структурного несоответствия были выявлены дополнительные ограничения:
% \begin{itemize}
%     \item \textit{Семантическая разреженность.} Во многих доменных корпусах (например, рецепты) набор возможных отношений крайне узок и не позволяет проверять устойчивость типизации для разнообразных классов связей.
%     \item \textit{Фрагментация при субсэмплинге.} Реальные пользовательские хранилища часто велики по объему; при выделении подмножества возникает высокий риск ``повисших'' ссылок (dangling links) и потери контекста, что смещает оценочные метрики.
%     \item \textit{Зашумленность текста.} Обрывочные формулировки, телеграфный стиль, вставки кода и цитаты без контекста усложняют извлечение сущностей и повышают неопределенность в решениях моделей.
%     \item \textit{Лицензирование.} Во многих персональных репозиториях отсутствует явная лицензия либо присутствуют фрагменты, потенциально нарушающие авторские права, что делает публикацию производных размеченных данных рискованной.
% \end{itemize}

% \textbf{Выбор Википедии как базового источника корпуса.}
% Для выполнения требований к атомарности и воспроизводимости 
% был выбран корпус статей Википедии, преобразованный в формат Obsidian-хранилища. 
% Такой выбор обоснован следующими свойствами:
% \begin{itemize}
%     \item \textit{Естественная атомарность.} Статья энциклопедии в большинстве случаев соответствует одной сущности/концепту, что согласуется с принципом ``одна заметка = одна сущность''.
%     \item \textit{Наличие базовой топологии.} Гиперссылки между статьями образуют связный граф, который удобен как источник кандидатных ребер и как слабая разметка для последующей валидации.
%     \item \textit{Качество текста.} Энциклопедический стиль обеспечивает высокую информационную плотность и связность изложения, что благоприятно для извлечения контекстов и принятия решений LLM.
%     \item \textit{Воспроизводимость сбора.} Сбор может быть формализован через публичный API, фиксируемые параметры обхода и детерминированные правила очистки/нормализации.
% \end{itemize}

% \textbf{Построение Markdown-корпуса и восстановление связей.}
% Подготовка корпуса выполнялась программно на основе Wikipedia API с использованием библиотеки \texttt{wikipediaapi} и включала следующие шаги.
% \begin{enumerate}
%     \item \textit{Выбор подграфа статей.} Реализованы два режима:
%           (i) целевая выгрузка заданного набора заголовков статей;
%           (ii) обход ``hub-and-spoke'' от одной стартовой темы (\texttt{seed}) с ограничением глубины и числа страниц.
%           Во втором режиме исходящие ссылки ранжируются кросс-энкодером (reranker), что позволяет сохранять более релевантные теме статьи при фиксированном бюджете на число заметок.
%     \item \textit{Нормализация идентификаторов.} Имена файлов приводятся к стабильному формату (замена пробелов на подчеркивания, удаление некоторых служебных символов), чтобы wikilink'и однозначно соответствовали путям файлов.
%     \item \textit{Очистка текста.} Из текста удаляются маркеры сносок (например, \texttt{[1]}), нормализуются подряд идущие переводы строк. Дополнительно из конца статей отсекаются секции, не несущие содержательной информации для извлечения отношений (например, ``См. также'', ``Литература'', ``Ссылки'' и их английские аналоги).
%           % \item \textit{Формирование блока связей.} Для каждой сохраненной статьи формируется секция \texttt{## Related Connections} со списком исходящих wikilink'ов. Для предотвращения ``повисших'' ребер ссылки на статьи, отсутствующие в выбранном подмножестве, удаляются.
%     \item \textit{Фильтрация ссылок по тексту (опционально).} Для снижения доли навигационных и слабосодержательных ссылок предусмотрен режим, в котором исходящая ссылка сохраняется только при наличии упоминания целевой сущности в основном тексте статьи (на основе лемматизации и проверки близости словоформ в предложениях).
% \end{enumerate}

% \textbf{Формирование ``золотого стандарта'' для оценки.}
% Для оценки качества системы требуются данные в двух представлениях.
% \begin{itemize}
%     \item \textit{Граф как Obsidian-хранилище.} Эталонные ребра фиксируются непосредственно в Markdown-файлах в виде типизированных строк внутри секции связей: \texttt{- relation:: [[Target]]}. Такая форма разметки позволяет (i) сохранять данные в естественном для Obsidian виде и (ii) напрямую сравнивать предсказания системы с эталоном путем парсинга Markdown.
%     \item \textit{Датасет контекстов упоминаний для retrieval-оценки.} Для каждой связи ``источник~--- цель'' извлекаются предложения контекста, в которых в тексте источника упоминается целевая сущность. Данные сериализуются в JSON и включают поля \texttt{source\_file}, \texttt{entity}, \texttt{sentence}, \texttt{is\_direct\_mention} и (после добавления типизации) \texttt{link\_type}.
% \end{itemize}
% Такое раздельное представление позволяет оценивать как задачу \textit{link classification} (типизация ребер), так и качество выделения релевантных контекстов, необходимых для LLM-классификатора и для GraphRAG-поиска.

% \textbf{Датасет и протокол оценки retrieval.}
% Оценка retrieval-компонента формализуется как задача воспроизведения контекстов упоминаний: для каждой пары \((s, e)\), где \(s\)~--- файл-источник (\texttt{source\_file}), а \(e\)~--- целевая сущность (\texttt{entity}), система должна вернуть хотя бы одно предложение/фрагмент текста (\texttt{sentence}), семантически совпадающий с эталонным контекстом из ``золотого стандарта''.
% Эталонные данные хранятся в виде набора JSON-файлов (по одному на подкорпус) в директории \texttt{data/test\_retrieval/gold/}. Каждый элемент массива соответствует одной ``истинной'' паре \((s, e)\) и содержит конкретное предложение из текста источника, в котором упоминается сущность \(e\).

% Предсказания retrieval-стратегий вычисляются непосредственно на Markdown-хранилищах подкорпусов (директория \texttt{data/test\_vaults/gold/}) и сохраняются в том же JSON-формате. В экспериментальном протоколе рассматриваются три стратегии:
% \begin{itemize}
%     \item \textit{Strict}~--- извлечение кандидатных упоминаний на основе лемматизированного совпадения сущностей в тексте;
%     \item \textit{Broad}~--- извлечение кандидатных упоминаний через векторный поиск по чанкам с последующим переранжированием кросс-энкодером;
%     \item \textit{Combined}~--- объединение \textit{Strict} и \textit{Broad} с устранением дубликатов.
% \end{itemize}

% Сопоставление эталона и предсказаний производится не по точному совпадению строк, а по ``нечеткому'' критерию на уровне лемм. Для каждого эталонного предложения вычисляется множество лемм (без стоп-слов и пунктуации) и проверяется, что доля совпавших лемм между эталоном и предсказанным фрагментом превышает фиксированный порог. Такая процедура делает метрику устойчивой к незначительным различиям в пунктуации и форматировании и к тому, что retrieval может возвращать фрагмент, покрывающий эталонное предложение не дословно.
% Итоговая оценка агрегируется через микро-метрики \(Precision\), \(Recall\) и \(F1\) по счетчикам \(TP/FP/FN\), где:
% \begin{itemize}
%     \item \(TP\): эталонный контекст для пары \((s,e)\) найден (существует предсказанный фрагмент, прошедший ``нечеткое'' сопоставление);
%     \item \(FN\): эталонный контекст отсутствует среди предсказаний;
%     \item \(FP\): предсказана пара \((s,e)\), для которой не удалось сопоставить ни один предсказанный фрагмент с эталонными предложениями.
% \end{itemize}

% \textbf{Датасет и протокол оценки LLM-классификации отношений.}
% Оценка LLM-компонента трактуется как многоклассовая классификация типа отношения для пар \((s,e)\) при условии, что retrieval уже предоставил набор релевантных контекстов. В качестве ``золотого стандарта'' используется тот же набор JSON-файлов \texttt{data/test\_retrieval/gold/}, но дополненный меткой \texttt{link\_type}, которая соответствует типизированной ссылке в Markdown-хранилище и интерпретируется как истинный класс (\texttt{true\_type}).

% Для каждой пары \((s,e)\) все доступные предложения-контексты агрегируются и далее подаются в модуль ранжирования контекстов. Ранжирование выполняется кросс-энкодером (модель переранжирования из конфигурации проекта), который получает на вход запрос вида ``What is the relationship between \dots?'' и набор кандидатных предложений; после чего выбираются \textit{top-n} наиболее релевантных предложений. При ранжировании дополнительно приоритизируются предложения, отмеченные признаком прямого упоминания (\texttt{is\_direct\_mention}).

% Существенной особенностью протокола является то, что пространство классов задается не глобально, а на уровне каждого подкорпуса. Для каждого sample список допустимых типов связей \texttt{llm\_link\_types} извлекается из файла конфигурации \texttt{.kg\_builder/config.json}, что позволяет:
% \begin{itemize}
%     \item ограничивать выбор LLM только релевантными типами для конкретной предметной области;
%     \item исключать классы, отсутствующие в эталонной разметке данного поднабора;
%     \item корректно сравнивать результаты между подкорпусами, не смешивая различные онтологии связей.
% \end{itemize}

% Запросы к LLM выполняются батчами; после каждого батча результаты сохраняются в JSONL-файл-чекпоинт, что обеспечивает воспроизводимость и возможность продолжить оценку после прерывания. Для качества классификации вычисляются precision/recall/F1 по каждому классу, а также macro- и weighted-усреднения (веса пропорциональны поддержке класса в эталоне). Метрики рассчитываются отдельно по каждому подкорпусу и в совокупности по всем подкорпусам.

% \textbf{Контрольные подкорпуса.}
% Для проверки устойчивости алгоритма к различным видам структур и сложностей были сформированы несколько небольших ``контрольных'' хранилищ, каждое из которых соответствует отдельному поднабору (sample). Подкорпуса различаются тематикой, плотностью ссылок и специфическими ``сложными случаями'' (например, омонимия и конкурирующие смыслы сущностей).
% Сводная статистика по использованным подкорпусам приведена в таблице~\ref{tab:dataset_samples}.

% \begin{table}[H]
%     \centering
%     \caption{Сводная статистика контрольных хранилищ (``золотой стандарт'')}
%     \label{tab:dataset_samples}
%     \begin{tabular}{lccc}
%         \hline
%         Подкорпус              & Число заметок & Число типизированных ссылок & Число контекстных пар \\
%         \hline
%         \texttt{hub\_n\_spoke} & 15            & 77                          & 96                    \\
%         \texttt{realistic}     & 15            & 79                          & 103                   \\
%         \texttt{versus}        & 15            & 44                          & 62                    \\
%         \texttt{sequence}      & 7             & 23                          & 44                    \\
%         \texttt{personal}      & 11            & 23                          & 22                    \\
%         \hline
%     \end{tabular}
% \end{table}

% Итоговый выбор в пользу Википедии и построенных на ее основе контрольных хранилищ обеспечил выполнение ключевого методологического требования атомарности сущностей и позволил получить воспроизводимый набор данных с управляемым объемом разметки, достаточный для количественной оценки качества GraphRAG-конвейера.


\clearpage
\section{Эксперименты}
\subsection{Параметры воспроизводимости}

Для обеспечения воспроизводимости экспериментов
был зафиксирован сид генерации (\texttt{seed})
на значении, равном 42.
Также во всех экспериментах с LLM
(как при типизации связей, так и при
генерации ответов)
температура генерации ($\text{temperature}$)
равнялась $0$.

% % ============================================================
\subsection{Оценка конвейера построения графа знаний}
\subsubsection{Извлечение кандидатных пар}
% % ============================================================


Извлечение может быть осуществлено тремя способами: \texttt{Strict}, \texttt{Broad} и \texttt{Combined}.
Подробное описание этих стратегий приведено в разделе 2.1.

Описание методологии оценивания изложено в разделе 1.4 «Оценка извлечения контекста».
Скрипт оценивания расположен в директории \texttt{src/eval/retrieval}.

\subsubsection*{Стратегия Strict}
В таблице~\ref{tab:strict_retrieval} приведены результаты работы стратегии
\texttt{Strict} на валидационном датасете для алгоритмов извлечения.

\begin{table}[H]
    \caption{Результат \texttt{Strict} на валидационном датасете}
    \label{tab:strict_retrieval}
    \centering
    \begin{tabular}{cccc}
        \hline
        \textbf{Precision} & \textbf{Recall} & \textbf{F\textsubscript{1}} & \textbf{Время, с} \\
        \hline
        1.0                & 0.376           & 0.547                       & 206               \\
        \hline
    \end{tabular}
\end{table}

Список извлеченных фрагментов не содержит ложноположительных результатов,
что обеспечивает идеальную точность. Однако такой подход ограничен по полноте,
так как он не способен извлечь чанки, в которых искомая сущность не упомянута явно.
Таким образом, подход задает верхнюю границу по точности и
нижнюю по полноте.
Стоит отметить, что эта стратегия извлекла все прямые упоминания сущностей.

Общее время работы стратегии на всем датасете составило 206 секунд,
что составляет около 41 секунды на подкорпус в среднем.
Кроме того, лемматизация текста заняла около 194 секунд,
а сам поиск~--- оставшиеся 12 секунд.

\subsubsection*{Стратегия Broad}

Качество работы данной стратегии зависит от значений параметров, а также используемых моделей,
поэтому для достижения наилучшего качества извлечения был проведен ряд экспериментов.

Первоначальные настройки указаны в таблице~\ref{tab:initial_hyperparameters}.
Их выбор был основан на значениях, используемых в референсных решениях
ObsidianRAG~\cite{obsidianrag} и Neo4j Graph Builder~\cite{neo4j_graph_builder},
а также на структуре валидационного датасета.
Поскольку в нем используются markdown-файлы,
имеет смысл приоритизировать разбиение текста на основе
заголовков.
Модели выбраны на основе их метрик и популярности на HuggingFace.

Инференс Embedding и Cross Encoder моделей осуществлялся локально
под управлением операционной системы Arch Linux,
оснащенной графическим ускорителем AMD Radeon RX 6700 XT с 12 ГБ видеопамяти.
Аппаратное ускорение обеспечивалось платформой ROCm (Radeon Open Compute).
Все модели использовались в формате \texttt{bfloat16},
что соответствует их оригинальной спецификации.

\begin{table}[H]
    \caption{Изначальные гиперпараметры алгоритма извлечения}
    \label{tab:initial_hyperparameters}
    \centering
    \renewcommand{\arraystretch}{1.0}
    \begin{tabularx}{\textwidth}{|>{\raggedright\arraybackslash}p{6cm}|>{\raggedright\arraybackslash}X|}
        \hline
        \textbf{Параметр}     & \textbf{Значение}                                                                                                                                                             \\
        \hline
        \multicolumn{2}{|c|}{\textbf{Эмбеддинги и разбиение текста}}                                                                                                                                          \\
        \hline
        Модель эмбеддингов    & \texttt{paraphrase-multilingual-MiniLM-L12-v2~\cite{minilm_v2}}                                                                                                               \\
        \hline
        Размер чанка          & 4000                                                                                                                                                                          \\
        \hline
        Перекрытие чанков     & 200                                                                                                                                                                           \\
        \hline
        Разделители           & \texttt{["\textbackslash{}n\#~", "\textbackslash{}n\#\#~", "\textbackslash{}n\#\#\#~", "\textbackslash{}n\textbackslash{}n", ".~", "?~", "!~", "\textbackslash{}n", "~", ""]} \\
        \hline
        \multicolumn{2}{|c|}{\textbf{Поиск}}                                                                                                                                                                  \\
        \hline
        Top-K                 & 15                                                                                                                                                                            \\
        \hline
        Вес векторного поиска & 0.5                                                                                                                                                                           \\
        \hline
        \multicolumn{2}{|c|}{\textbf{Переранжирование}}                                                                                                                                                       \\
        \hline
        Модель                & \texttt{BAAI/bge-reranker-v2-m3~\cite{bge_reranker_v2}}                                                                                                                       \\
        \hline
        Top-K                 & 5                                                                                                                                                                             \\
        \hline
        Порог                 & 0.7                                                                                                                                                                           \\
        \hline
    \end{tabularx}
\end{table}

Для поиска оптимальной конфигурации был проведен ряд экспериментов,
заключающихся в смене размера чанка, моделей эмбеддингов и переранжирования.
Их результаты представлены в
таблице~\ref{tab:broad_rest}.

Поскольку стратегия \texttt{Strict}
уже обеспечивает идеальную точность,
основной идеей экспериментов над \texttt{Broad}
было увеличение полноты при
сохранении приемлемой точности.
Допустимый порог Precision был сделан равным 0.5.
Его достижение будет первостепенной задачей,
после чего приоритет сместится на максимизацию полноты
при сохранении точности выше порога.

\begin{table}[H]
    \caption{Анализ конфигураций стратегии \texttt{Broad}.}
    \label{tab:broad_rest}
    \centering
    \begin{tabular}{ccclclccccc}
        \hline
        \textbf{\#}                         & \textbf{CS} & \textbf{CO} & \textbf{E} & \textbf{K} & \textbf{CE} & \textbf{Порог} & \textbf{P} & \textbf{R} & \textbf{F\textsubscript{1}} & \textbf{Время} \\
        \hline
        1                                   & 4000        & 200         & MiniLM     & 15         & bge-r       & 0.7            & 0.354      & 0.199      & 0.255                       & 34             \\
        2                                   & 3000        & 200         & MiniLM     & 15         & bge-r       & 0.7            & 0.395      & 0.190      & 0.256                       & 40             \\
        3                                   & 5000        & 300         & MiniLM     & 15         & bge-r       & 0.7            & 0.301      & 0.186      & 0.223                       & 27             \\
        4                                   & 3000        & 200         & MiniLM     & 15         & jina        & 0.7            & 0.553      & 0.208      & 0.302                       & 9              \\
        5                                   & 3000        & 200         & MiniLM     & 15         & jina        & 0.5            & 0.510      & 0.407      & 0.452                       & 9              \\
        6                                   & 3000        & 200         & MiniLM     & 30         & jina        & 0.5            & 0.500      & 0.453      & 0.475                       & 20             \\
        7                                   & 3000        & 200         & bge-e      & 30         & jina        & 0.5            & 0.531      & 0.468      & 0.498                       & 23             \\
        \hline
        \multicolumn{4}{l}{\textit{Strict}} &             &             &            & 1.0        & 0.376       & 0.547          & 3.4                                                                    \\
        \hline
    \end{tabular}

    \vspace{1ex}
    \raggedright\footnotesize
    \textit{Примечание:} Используемые модели: \\
    MiniLM~--- \texttt{paraphrase-multilingual-MiniLM-L12-v2}~\cite{minilm_v2}, \\
    bge-e~--- \texttt{BAAI/bge-m3}~\cite{bge_m3}, \\
    bge-r~--- \texttt{BAAI/bge-reranker-v2-m3}~\cite{bge_reranker_v2}, \\
    jina~--- \texttt{jina-reranker-v2-base-multilingual}~\cite{jina_reranker_v2}. \\
    CS~--- Chunk Size, CO~--- Chunk Overlap,
    E~--- Embedding Model, CE~--- Cross Encoder, Порог~--- минимально допустимая оценка CE,
    время в минутах
\end{table}

\subsubsection*{Влияние размера чанка}
Эксперименты 1-3 показали, что фрагмент текста не может быть слишком большим,
так как это приводит к потере деталей, что негативно сказывается на точности поиска.
Оптимальным по качеству вариантом оказался размер в 3000 символов с перекрытием в 200 символов.

Стоит отметить влияние размера чанков на время работы.
Уменьшение размера чанка ведет к увеличению числа
фрагментов, которые нужно перевести в векторное пространство,
из-за чего растет количество пар для сравнения,
и следовательно время работы.
Поэтому возникает компромисс:
увеличить размер чанков, сделав их менее детализированными, но снизив время работы,
либо уменьшить размер, увеличив время.

\subsubsection*{Роль модели переранжирования}
Модель Cross Encoder также оказалась крайне важной.
Переход с bge-r на jina увеличил $\text{F}_1$ до 0.302 (эксперимент 4),
но потребовал снижения порога с 0.7 до 0.5 для раскрытия потенциала модели (эксперимент 5).
Как итог, $\text{F}_1$ вырос до 0.452.

Изначальное значение порога оценки Cross Encoder для фильтрации результатов
было завышенным, из-за чего
часть истинно положительных кандидатов отбрасывалась.
При bge-r порог в 0.7 был оправдан, так как
модель имела низкую точность, но использование jina
решило эту проблему, повысив значение Precision с 0.395 до 0.553.

Смена модели также значительно сократила время работы алгоритма
(с 40 до 9 минут). Это объясняется вдвое меньшим количеством параметров
в модели jina (300 миллионов против 600 у bge-r),
а также более современной и вычислительно эффективной архитектурой.

Так как первостепенная цель экспериментов достигнута (порог превысил 0.5),
последующие опыты будут нацелены на максимизацию полноты
при сохранении приемлемой точности.

\subsubsection*{Количество извлекаемых чанков}
С ростом точности извлечения и значительным сокращением
времени работы появился смысл увеличить число
извлекаемых кандидатов при гибридном поиске (Top-K).
Снижение его значения не имеет особого смысла,
так как это скорее всего сильно снизит полноту,
что противоречит цели.

По результатам эксперимента 6 видно,
что увеличение Top-K спровоцировало рост шума,
а также времени работы, но увеличило полноту.

\subsubsection*{Смена модели эмбеддингов}
На замену MiniLM, переводящую тексты в 384-мерное пространство
пришла более тяжелая bge-m3, которая работает с 1024-мерными эмбеддингами.
Идея была в улучшении сохранения детализации смыслов в текстах.
Как итог, рост точности и полноты на 0.031 и 0.015 соответственно
по сравнению с результатами эксперимента 6.

\subsubsection*{Стратегия Combined}

В данной стратегии использовалась оптимальная конфигурация \texttt{Broad}.
Результаты представлены в таблице~\ref{tab:combined_retrieval}.
Видно, что \texttt{Broad} вносит весомый вклад:
полнота выросла практически вдвое (с 0.376 до 0.685),
но ожидаемо упала точность (с 1 до 0.675).
Полученные результаты подтверждают
теорию: данная стратегия действительно является компромиссом
между высокой точностью и полнотой.

\begin{table}[H]
    \caption{Результат \texttt{Combined} на валидационном датасете}
    \label{tab:combined_retrieval}
    \centering
    \begin{tabular}{clllccc}
        \hline
        \textbf{Precision} & \textbf{Recall} & \textbf{F\textsubscript{1}} \\
        \hline
        0.675              & 0.685           & 0.680                       \\
        \hline
    \end{tabular}
\end{table}

\begin{figure}[H]
    \centering
    \begin{tikzpicture}
        \begin{axis}[
                width=0.85\textwidth,
                height=6cm,
                xlabel={Номер эксперимента},
                ylabel={F\textsubscript{1}},
                xmin=0.5, xmax=7.5,
                ymin=0, ymax=0.80,
                xtick={1,2,3,4,5,6,7},
                ytick={0,0.1,0.2,0.3,0.4,0.5,0.6,0.7,0.8},
                ymajorgrids=true,
                grid style=dashed,
                legend style={at={(0.02,0.98)}, anchor=north west, font=\small},
                tick label style={font=\small},
                label style={font=\small},
            ]
            \addplot[color=blue, mark=*, thick] coordinates {
                    (1, 0.255)
                    (2, 0.256)
                    (3, 0.223)
                    (4, 0.302)
                    (5, 0.452)
                    (6, 0.475)
                    (7, 0.498)
                };
            \addlegendentry{Broad}

            \addplot[color=red, dashed, thick, mark=none] coordinates {
                    (0.5, 0.547) (7.5, 0.547)
                };
            \addlegendentry{Strict (F\textsubscript{1}=0.547)}

            \addplot[color=teal, only marks, mark=square*, mark size=4pt] coordinates {
                    (7, 0.680)
                };
            \addlegendentry{Combined (эксп. 7)}

        \end{axis}
    \end{tikzpicture}
    \caption{Динамика F\textsubscript{1} стратегии \texttt{Broad} по экспериментам.
        Красная пунктирная линия~--- уровень \texttt{Strict}-стратегии.
        Квадрат~--- результат \texttt{Combined} для лучшей конфигурации.}
    \label{fig:retrieval_f1}
\end{figure}

% % ============================================================
\subsubsection{Типизация связей}
\label{sec:llm_eval}
% % ============================================================

В рамках данной серии экспериментов
рассматривалась способность больших языковых моделей
корректно определять семантический тип связи между сущностями
на основе собранных контекстов.
Для оценки использовался вручную размеченный валидационный датасет.

Поскольку важна как точность, так и полнота классификации,
в качестве основной метрики был выбран F\textsubscript{1}-score.
Кроме того, так как в данных присутствует дисбаланс,
уместнее всего использовать Weighted F\textsubscript{1}.
Идея экспериментов заключалась в максимизации основной метрики
путем поиска оптимального промпта и модели.

Результаты опытов представлены в таблице~\ref{tab:classification_results}.
LLM модели использовались посредством API облачных платформ.
Для Qwen-3-235B~\cite{qwen_3_235b} использовался Cerebras Inference API~\cite{cerebras_api},
для Gemma-3-27B~\cite{gemma_3_27b} и Gemma-4-31B~\cite{gemma_4_31b}~--- Google Vertex AI~\cite{vertex_ai}.
Оценивание производилось согласно методологии,
описанной в разделе 1.4 «Оценка LLM-типизации».
Скрипт оценивания расположен в директории \texttt{src/eval/classification}.

\begin{table}[H]
    \caption{Weighted F\textsubscript{1} классификации типов связей}
    \label{tab:classification_results}
    \centering
    \begin{tabular}{llccccccc}
        \hline
        \textbf{Модель} & \textbf{Промпт} & \textbf{Batch}                 &
        \texttt{P}      & \texttt{H}      & \texttt{R}                     &
        \texttt{S}      & \texttt{V}      & \textbf{W. F\textsubscript{1}}                                                          \\
        \hline
        Gemma-3-27B     & base            & 1                              & 0.346 & 0.340 & 0.351 & 0.459 & 0.522 & 0.393          \\
        Gemma-3-27B     & base            & 5                              & 0.446 & 0.325 & 0.321 & 0.451 & 0.573 & 0.387          \\
        Gemma-3-27B     & base            & 10                             & 0.378 & 0.428 & 0.363 & 0.392 & 0.572 & 0.423          \\
        Qwen-3-235B     & base            & 1                              & 0.313 & 0.454 & 0.272 & 0.454 & 0.526 & 0.392          \\
        Gemma-3-27B     & few-shot        & 1                              & 0.372 & 0.436 & 0.444 & 0.341 & 0.559 & 0.445          \\
        Qwen-3-235B     & few-shot        & 1                              & 0.498 & 0.390 & 0.434 & 0.558 & 0.482 & 0.445          \\
        Gemma-4-31B     & few-shot        & 1                              & 0.856 & 0.578 & 0.627 & 0.786 & 0.758 & \textbf{0.669} \\
        \hline
    \end{tabular}

    \vspace{1ex}
    \raggedright\footnotesize
    \textit{Примечание:} Используемые промпты
    доступны в репозитории проекта по пути: \texttt{artifacts/prompts/classification.} \\
\end{table}

\subsubsection*{Влияние пакетного промптинга}
Увеличение количества пар заметок, передаваемых модели в рамках одного запроса,
в среднем снижает среднюю оценку, однако в некоторых случаях она растет.
Увеличение может быть вызвано утечкой контекста:
модель получает дополнительные данные из других элементов пакета.
Это видно в случае Personal, Hub-n-Spoke и Versus, где размер содержимого заметок мал,
и есть кластеры из тематически связных заметок.
Тем не менее, важно не перегрузить модель лишним контекстом: при классификации конкретной связи
контекст для других считается шумом и не должен быть учтен. Не все модели могут корректно
разграничить содержимое, из-за чего начинают путаться и делать ошибочные выводы, ухудшая качество классификации.
Как наиболее честный вариант был выбран размер пакета равный 1.

\subsubsection*{Влияние промпта}
Добавление в промпт примеров и описаний типов связей, которые в них используются,
увеличило значение Overall F\textsubscript{1} как для Gemma-3-27B с 0.392 до 0.423 ~(+0.031),
так и для Qwen-3-235B с 0.392 до 0.445~(+0.135).
Это подтверждает чувствительность моделей
к инструкциям и примерам: они помогают им лучше понять специфику задачи
и не уйти в сторону при рассуждении.
Важно точно определить каждый тип связи для устранения неоднозначностей.
Кроме того, результаты показывают,
что качество более крупной модели больше зависит от промпта.

\subsubsection*{Влияние выбора модели}
Эксперименты показали, что
эффективность архитектуры важнее числа параметров.
Новая Gemma-4-31B (31 млрд параметров) значительно превосходит
по качеству Qwen-3-235B (235 млрд параметров),
а Gemma-3-27B (27 млрд параметров) оказалась в ней наравне.

\subsubsection*{Анализ ошибок классификации}
Несмотря на рост общего значения Weighted F\textsubscript{1},
имеется несколько общих ошибок:
\begin{itemize}
    \item Злоупотребление универсальными типами.
          Тип \texttt{mentions} является самым популярным в датасете
          (40 связей, F\textsubscript{1} = 0.23--0.39).
          Он играет роль резервного класса:
          он назначается в случае,
          когда связь есть, но невозможно однозначно назначить другой класс.
          Модель может слишком часто назначать этот тип из-за его универсальности.
          Эта проблема решается дополнительными инструкциями и примерами в промпте,
          явно указывающими смысл этой связи.
    \item Использование близких по смыслу связей.
          Классы \texttt{encompasses}, \texttt{contains}, \texttt{incorporates} и другие
          имеют высокое семантическое сходство.
          Большое их количество создает путаницу:
          модель не способна выявить семантической разницы между ними,
          поэтому с высокой вероятностью может использовать неверный тип.
          Проблема обостряется в реальных условиях неточного контекста,
          который создает алгоритм извлечения контекста.
          Эффективным решением является объединение
          подобных классов в один более универсальный.
\end{itemize}

\subsection{Сравнение сгенерированных ссылок с эталонными}

Перед оценкой GraphRAG на сгенерированных ссылках необходимо их проанализировать.
Оценка проводилась в двух режимах: без учета типа (оценивается только факт наличия связи между заметками)
и с учетом типа (оценивается как наличие связи, так и правильность предсказанного типа)
для анализа качества ссылок для использования в каждом типе обхода.
Использовались лучшие конфигурации извлечения (стратегия \texttt{Combined}) и типизации
(Gemma-4-31B, few-shot промпт, \texttt{batch=1}).

\begin{table}[H]
    \caption{Оценка качества сгенерированных связей: факт наличия (Untyped) и точное совпадение типа (Typed)}
    \label{tab:generated_links_eval}
    \centering
    \begin{tabular}{lcccccc}
        \hline
        \multirow{2}{*}{\textbf{Подкорпус}} & \multicolumn{3}{c}{\textbf{Без учета типа (Untyped)}} & \multicolumn{3}{c}{\textbf{С учетом типа (Typed)}}                                                                                       \\
        \cline{2-7}
                                            & \textbf{P}                                            & \textbf{R}                                         & \textbf{F\textsubscript{1}} & \textbf{P} & \textbf{R} & \textbf{F\textsubscript{1}} \\
        \hline
        \texttt{hub\_n\_spoke}              & 0.810                                                 & 0.883                                              & 0.845                       & 0.643      & 0.701      & 0.671                       \\
        \texttt{personal}                   & 0.909                                                 & 0.435                                              & 0.588                       & 0.909      & 0.435      & 0.588                       \\
        \texttt{realistic}                  & 0.771                                                 & 0.810                                              & 0.790                       & 0.687      & 0.722      & 0.704                       \\
        \texttt{sequence}                   & 0.733                                                 & 0.957                                              & 0.830                       & 0.533      & 0.696      & 0.604                       \\
        \texttt{versus}                     & 0.727                                                 & 0.909                                              & 0.808                       & 0.545      & 0.682      & 0.606                       \\
        \hline
        \textbf{Взвеш. усред.}              & 0.776                                                 & 0.829                                              & 0.802                       & 0.635      & 0.679      & 0.656                       \\
        \hline
    \end{tabular}
\end{table}

Точность извлеченных ссылок без учета их типа находится на высоком уровне:
$\text{F}_1 = 0.802$, при этом полнота также имеет близкое к идеальному значение
$\text{R} = 0.829$, что говорит о способности алгоритма находить
практически все ссылки.
Наилучшие результаты наблюдаются в подкорпусах
\texttt{sequence} ($\text{R} = 0.957$) и \texttt{versus} ($\text{R} = 0.909$).

Введение типов ссылок как критерия оценивания
снизило значения метрик: $\text{F}_1$ упал до $0.656$.
Исключением является образец \texttt{personal}, в котором метрики не изменились.
Это означает, что алгоритм безошибочно классифицировал все связи.
Такой результат был ожидаем, он напрямую коррелирует с результатами раздела~\ref{sec:llm_eval}.

% % ============================================================
\subsection{Оценка конвейера GraphRAG}
\label{sec:graphrag_eval}
% % ============================================================

В таблице~\ref{tab:graphrag_default_params} представлены первоначальные параметры алгоритмов.
Они выбраны на основе количества валидационных данных и референсного решения.
Значения подобраны так, чтобы заставить алгоритмы обходить граф,
а не собрать контекст всех заметок разом.

\begin{table}[H]
    \caption{Значения гиперпараметров алгоритма GraphRAG по умолчанию}
    \label{tab:graphrag_default_params}
    \centering
    \renewcommand{\arraystretch}{1.0}
    \begin{tabularx}{\textwidth}{|>{\raggedright\arraybackslash}X|l|c|}
        \hline
        \textbf{Описание параметра}            & \textbf{Переменная}       & \textbf{Значение} \\
        \hline
        Максимальная глубина обхода графа      & \texttt{MAX\_HOPS}        & 2                 \\
        \hline
        Порог отсечения узлов по релевантности & \texttt{SCORE\_THRESHOLD} & 0.2               \\
        \hline
        Количество первоначальных заметок      & \texttt{TOP\_K\_SEED}     & 3                 \\
        \hline
        Итоговое количество чанков в контексте & \texttt{TOP\_K\_CONTEXT}  & 10                \\
        \hline
    \end{tabularx}
\end{table}

Оценка проводилась в соответствии с методологией,
описанной в разделе 1.4 «Полная оценка GraphRAG».
Скрипты оценивания расположены в директории \texttt{src/eval/ragas}.

\subsubsection{Анализ референсных решений}
\label{sec:reference_results}

Для оценки разработанных алгоритмов было проведено
их комплексное сравнение с ObsidianRAG и Neo4j
(Graph Builder и его же интегрированная реализация RAG).
В качестве модели для генерации ответов была выбрана Gemma-3-27B,
для оценки Gemma-4-31B.

\begin{table}[H]
    \caption{Результаты решения ObsidianRAG}
    \label{tab:reference_results_obsidianrag}
    \centering
    \begin{tabular}{lcccc}
        \hline
        \textbf{Подкорпус}     & \textbf{CR} & \textbf{CP} & \textbf{F} & \textbf{AC} \\
        \hline
        \texttt{hub\_n\_spoke} & 0.802       & 0.701       & 0.955      & 0.542       \\
        \texttt{personal}      & 0.625       & 0.615       & 1.000      & 0.424       \\
        \texttt{realistic}     & 0.886       & 0.707       & 0.898      & 0.505       \\
        \texttt{sequence}      & 1.000       & 0.824       & 0.967      & 0.582       \\
        \texttt{versus}        & 0.617       & 0.467       & 0.883      & 0.428       \\
        \hline
        \textbf{Среднее}       & 0.786       & 0.663       & 0.941      & 0.496       \\
        \hline
    \end{tabular}
\end{table}

Результаты оценки референсного решения ObsidianRAG приведены в таблице~\ref{tab:reference_results_obsidianrag}.
Явно видны недостатки логики расширения контекста: в
\texttt{personal} и \texttt{versus} наблюдается низкий CR (0.625 и 0.617 соответственно).
В этих образцах есть два кластера, которые практически друг с другом не связаны.
Наивный сбор контекста соседних узлов не позволит получить информацию из другого кластера,
а лишь увеличит шум, что видно по низким значениям CP
(0.467 в случае versus и 0.615 в случае personal).
В других образцах, где заметки связаны большим числом ссылок, CR не ниже 0.8, а CP не ниже 0.7,
что говорит о хорошей работе алгоритма.

Как результат, следует и низкий AC:
контекст зашумлен и в нем недостаточно
данных для предоставления
полного и точного ответа.

\begin{table}[H]
    \caption{Результаты решения Neo4j на сгенерированном ГЗ}
    \label{tab:reference_results_neo4j}
    \centering
    \begin{tabular}{lcccc}
        \hline
        \textbf{Подкорпус}     & \textbf{CR} & \textbf{CP} & \textbf{F} & \textbf{AC} \\
        \hline
        \texttt{hub\_n\_spoke} & 1.000       & 0.782       & 0.951      & 0.598       \\
        \texttt{personal}      & 1.000       & 0.875       & 0.984      & 0.612       \\
        \texttt{realistic}     & 0.985       & 0.842       & 0.975      & 0.620       \\
        \texttt{sequence}      & 0.980       & 0.855       & 0.971      & 0.640       \\
        \texttt{versus}        & 0.920       & 0.751       & 0.905      & 0.540       \\
        \hline
        \textbf{Среднее}       & 0.977       & 0.821       & 0.957      & 0.602       \\
        \hline
    \end{tabular}
\end{table}

В таблице~\ref{tab:reference_results_neo4j} представлены результаты Neo4j,
полученные на сгенерированном Neo4j Graph Builder ГЗ.

Решение имеет значительно более высокие значения CR,
что обусловлено более продвинутой логикой сбора контекста.
Среднее значение CR достигло $0.977$,
в то время как у ObsidianRAG оно лишь $0.786$.
Также виден большой прирост метрики на подкорпусах
\texttt{personal} (с $0.625$ до $1.000$) и
\texttt{versus} (с $0.617$ до $0.920$).

Кроме того, CP также вырос~---
в среднем $0.821$ против $0.663$ у ObsidianRAG.
Это говорит об отличной способности Neo4j фильтровать контекст.

Как следствие, увеличились значения AC:
среднее значение выросло на $0.106$ балла (с $0.496$ до $0.602$).
Прирост особенно заметен на проблемных для ObsidianRAG подкорпусах
\texttt{personal} ($+0.188$) и \texttt{versus} ($+0.112$).

Во всех случаях наблюдается высокий Answer Faithfulness:
модель не пытается додумать ответ, а использует приведенные данные.

\subsubsection{Сравнение собственных алгоритмов с аналогами}

Поскольку основным преимуществом алгоритмов, представленных в данной работе,
является многошаговый обход, основная цель которого увеличение CR,
имеет смысл сфокусироваться только на полноте контекста и качестве ответа
при сравнении разработанных алгоритмов с эталонным.
Сравнение представлено в таблице~\ref{tab:complex_graphrag_comparison}.

\begin{table}[H]
    \caption{Сравнение собственных алгоритмов с аналогами на размеченном (GT) и сгенерированном (GEN) графах.}
    \label{tab:complex_graphrag_comparison}
    \centering
    \renewcommand{\arraystretch}{1.1}
    \setlength{\tabcolsep}{4pt}
    \begin{tabular}{l|l|cc|cc|cc|cc}
        \hline
        \multirow{2}{*}{\textbf{Подкорпус}} & \multirow{2}{*}{\textbf{Данные}} & \multicolumn{2}{c|}{\textbf{Ref}} & \multicolumn{2}{c|}{\textbf{Base}} & \multicolumn{2}{c|}{\textbf{Typed}} & \multicolumn{2}{c}{\textbf{Neo4j}}                                                         \\
        \cline{3-10}
                                            &                                  & \textbf{CR}                       & \textbf{AC}                        & \textbf{CR}                         & \textbf{AC}                        & \textbf{CR} & \textbf{AC} & \textbf{CR} & \textbf{AC} \\
        \hline
        \multirow{2}{*}{\texttt{hub}}
                                            & GT                               & 0.802                             & 0.542                              & 1.000                               & 0.587                              & 0.958       & 0.565       & ---         & ---         \\
                                            & GEN                              & 0.750                             & 0.516                              & 0.958                               & 0.535                              & 0.958       & 0.561       & 1.000       & 0.598       \\
        \hline
        \multirow{2}{*}{\texttt{personal}}
                                            & GT                               & 0.625                             & 0.424                              & 1.000                               & 0.597                              & 1.000       & 0.579       & ---         & ---         \\
                                            & GEN                              & 0.625                             & 0.429                              & 1.000                               & 0.602                              & 1.000       & 0.642       & 1.000       & 0.612       \\
        \hline
        \multirow{2}{*}{\texttt{realistic}}
                                            & GT                               & 0.886                             & 0.505                              & 0.977                               & 0.606                              & 0.977       & 0.601       & ---         & ---         \\
                                            & GEN                              & 0.864                             & 0.522                              & 0.977                               & 0.637                              & 0.977       & 0.602       & 0.985       & 0.620       \\
        \hline
        \multirow{2}{*}{\texttt{sequence}}
                                            & GT                               & 1.000                             & 0.582                              & 0.925                               & 0.580                              & 0.975       & 0.631       & ---         & ---         \\
                                            & GEN                              & 0.975                             & 0.555                              & 0.925                               & 0.614                              & 0.925       & 0.600       & 0.980       & 0.640       \\
        \hline
        \multirow{2}{*}{\texttt{versus}}
                                            & GT                               & 0.617                             & 0.428                              & 0.900                               & 0.524                              & 0.788       & 0.507       & ---         & ---         \\
                                            & GEN                              & 0.608                             & 0.388                              & 0.900                               & 0.550                              & 0.788       & 0.520       & 0.920       & 0.540       \\
        \hline
        \hline
        \multirow{2}{*}{\textbf{Среднее}}
                                            & GT                               & 0.786                             & 0.496                              & 0.960                               & 0.579                              & 0.940       & 0.577       & ---         & ---         \\
                                            & GEN                              & 0.764                             & 0.482                              & 0.952                               & 0.587                              & 0.930       & 0.585       & 0.977       & 0.602       \\
        \hline
    \end{tabular}
\end{table}

\subsubsection*{Оценка GraphRAG на размеченных данных}

Видно, что модифицированная логика обхода действительно
увеличила количество собираемой информации по сравнению с ObsidianRAG,
позволив получить практически идеальную полноту контекста.
Это оказалось особенно заметно в случае \texttt{personal} и \texttt{versus},
где изначально значение CR было самым низким.
В их случае прирост составил 0.375 и 0.283 соответственно в случае Base,
0.375 и 0.171 в случае Typed.

Дополнительный контекст помог языковой модели дать более точный ответ,
что также увеличило средние значение AC с
$0.496$ до $0.579$ (Base) и $0.577$ (Typed)

% % ============================================================
\subsubsection*{Оценка GraphRAG на сгенерированном графе}

Для получения первоначальных данных и оценки качества
работы алгоритмов на автоматически построенных ГЗ
в следующем эксперименте размеченный ГЗ был заменен
на граф, построенным пайплайном генерации связей.
Результаты приведены в таблице~\ref{tab:complex_graphrag_comparison}.

Видно, что замена графа минимально влияет на качество работы GraphRAG:
значения метрик колеблются в диапазоне сотых долей.

Переход на сгенерированные связи незначительно снизил CR:
для базового алгоритма среднее значение упало
с $0.960$ до $0.952$, а для типизированного~— с $0.940$ до $0.930$.
Это падение можно
объяснить пропуском неочевидных связей конвейером построения ГЗ,
из-за чего усложняется навигация по графу.

Интересным наблюдением является рост AC
(с $0.579$ до $0.587$ для Base и с $0.577$ до $0.585$ для Typed).
Особенно заметно увеличение в
подкорпусе \texttt{personal} для Typed-обхода (с $0.579$ до $0.642$)
и в \texttt{realistic} для Base-обхода (с $0.606$ до $0.637$).
Возможно, что сгенерированные связи в некоторых случаях
более эффективны для сбора релевантного контекста,
так как пайплайн мог иначе типизировать связи или создать новые,
которые помогли извлечь нужные данные.

Алгоритм Base показал себя лучше в большинстве случаев
как при сборе контекста, так и при генерации ответов.

Результаты работы предложенных алгоритмов оказались сопоставимы
с Neo4j: в среднем алгоритм Base уступает ему всего на $0.025$ по полноте и на $0.015$
по точности ответа, алгоритм Typed на $0.047$ и $0.017$  соответственно.
Это доказывает эффективность разработанных алгоритмов.

% ============================================================
\subsubsection{Влияние промпта и языковой модели}
\label{sec:llm_prompt_effect}

Для исследования влияния различных факторов и компонент GraphRAG на качество ответа LLM
была проведена серия экспериментов над решениями, представленными в данной работе.
Все замеры проводились на размеченных данных.

\subsubsection*{Базовый промпт против few-shot}
На основе экспериментов над LLM-классификацией связей был сделан вывод о важности
указания точных инструкций и конкретных примеров в промпте.
В базовом промпте уже есть основные указания,
такие как придерживание предоставленного контекста и ответ на языке запроса.
К ним были добавлены конкретные развернутые примеры и
требования модели к рассуждениям перед ответом.
Модели для оценки и генерации остались теми же.

\begin{table}[H]
    \caption{Сравнение базового и few-shot промптов на алгоритмах Base и Typed (метрика AC, модель Gemma-3-27B)}
    \label{tab:prompt_effect_full_versus}
    \centering
    \begin{tabular}{llcccccc}
        \hline
        \textbf{Алгоритм}      & \textbf{Промпт} & \texttt{H} & \texttt{P} & \texttt{R} & \texttt{S} & \texttt{V} & \textbf{AC avg} \\
        \hline
        \multirow{2}{*}{Base}  & Базовый         & 0.587      & 0.597      & 0.606      & 0.580      & 0.524      & 0.579           \\
                               & Few-shot        & 0.595      & 0.572      & 0.604      & 0.618      & 0.569      & 0.592           \\
        \hline
        \multirow{2}{*}{Typed} & Базовый         & 0.565      & 0.579      & 0.601      & 0.631      & 0.507      & 0.577           \\
                               & Few-shot        & 0.600      & 0.611      & 0.606      & 0.629      & 0.555      & 0.580           \\
        \hline
    \end{tabular}

    \vspace{1ex}
    \raggedright\footnotesize
    \textit{Примечание:} Используемые промпты
    доступны в репозитории проекта по пути: \texttt{artifacts/prompts/graphrag.} \\
\end{table}

Новый промпт продемонстрировал прирост AC в среднем лишь +0.13.
Прирост оказался ниже, чем в предыдущих сравнениях, вероятно,
из-за корректно сформулированных инструкций в базовом промпте.
Можно сделать вывод,
что наличие примеров все же дает небольшой прирост корректности ответа языковой модели.

\subsubsection*{Анализ влияния языковой модели}

В этом эксперименте
модель для генерации будет заменена на Gemma-4-31B, а
моделью-судьей будет Gemini-3-flash~\cite{gemini_3_flash}.

\begin{table}[H]
    \caption{Сравнение моделей Gemma-3-27B и Gemma-4-31B (few-shot промпт)}
    \label{tab:model_comparison_2_3}
    \centering
    \begin{tabular}{llccc|ccc}
        \hline
        \multirow{2}{*}{\textbf{Датасет}}   & \multirow{2}{*}{\textbf{Алгоритм}} & \multicolumn{3}{c|}{\textbf{Эксп. 2 (Gemma-3-27B)}} & \multicolumn{3}{c}{\textbf{Эксп. 3 (Gemma-4-31B)}}                                                                            \\
        \cline{3-8}
                                            &                                    & \textbf{CP}                                         & \textbf{CR}                                        & \textbf{AC} & \textbf{CP}        & \textbf{CR}        & \textbf{AC}      \\
        \hline
        \multirow{2}{*}{\texttt{hub}}       & Base                               & 0.749                                               & 1.000                                              & 0.595       & 0.660 $\downarrow$ & 0.875 $\downarrow$ & 0.698 $\uparrow$ \\
                                            & Typed                              & 0.796                                               & 0.958                                              & 0.600       & 0.707 $\downarrow$ & 0.833 $\downarrow$ & 0.709 $\uparrow$ \\
        \hline
        \multirow{2}{*}{\texttt{personal}}  & Base                               & 0.893                                               & 1.000                                              & 0.572       & 0.750 $\downarrow$ & 1.000 $=$          & 0.796 $\uparrow$ \\
                                            & Typed                              & 0.892                                               & 1.000                                              & 0.611       & 0.729 $\downarrow$ & 0.917 $\downarrow$ & 0.770 $\uparrow$ \\
        \hline
        \multirow{2}{*}{\texttt{realistic}} & Base                               & 0.848                                               & 0.977                                              & 0.604       & 0.426 $\downarrow$ & 1.000 $\uparrow$   & 0.738 $\uparrow$ \\
                                            & Typed                              & 0.851                                               & 0.977                                              & 0.606       & 0.491 $\downarrow$ & 0.955 $\downarrow$ & 0.792 $\uparrow$ \\
        \hline
        \multirow{2}{*}{\texttt{sequence}}  & Base                               & 0.879                                               & 0.925                                              & 0.618       & 0.434 $\downarrow$ & 0.765 $\downarrow$ & 0.734 $\uparrow$ \\
                                            & Typed                              & 0.878                                               & 0.975                                              & 0.629       & 0.429 $\downarrow$ & 0.765 $\downarrow$ & 0.706 $\uparrow$ \\
        \hline
        \multirow{2}{*}{\texttt{versus}}    & Base                               & 0.734                                               & 0.900                                              & 0.569       & 0.519 $\downarrow$ & 0.655 $\downarrow$ & 0.623 $\uparrow$ \\
                                            & Typed                              & 0.696                                               & 0.813                                              & 0.555       & 0.547 $\downarrow$ & 0.568 $\downarrow$ & 0.576 $\uparrow$ \\
        \hline
        \multirow{2}{*}{\textbf{Среднее}}   & Base                               & 0.821                                               & 0.960                                              & 0.592       & 0.558 $\downarrow$ & 0.859 $\downarrow$ & 0.718 $\uparrow$ \\
                                            & Typed                              & 0.823                                               & 0.945                                              & 0.600       & 0.581 $\downarrow$ & 0.808 $\downarrow$ & 0.711 $\uparrow$ \\
        \hline
    \end{tabular}
\end{table}

Результаты представлены в таблице~\ref{tab:model_comparison_2_3}.
Прирост AC составил в среднем +0.128,
при этом CP упал в среднем на 0.263.
Сильнее всего CP снизился на \texttt{realistic} (с 0.848 до 0.426)
и \texttt{sequence} (с 0.879 до 0.434).
Снижение CP обусловлено более строгой оценкой модели.
Однако рост AC показывает, что более мощная модель
для генерации ответов способна абстрагироваться от шума
и дать более точный ответ.

\subsubsection*{Сравнение типизированного обхода с обычным}

У типизированного обхода выше CP (среднее: 0.581 против 0.558 у Base), но ниже
CR (0.808 против 0.859 у Base).
Эти результаты показывают влияние фильтрации заметок при обходе графа:
тематически нерелевантные файлы отсекаются, но при этом также теряется возможность
сделать дополнительный шаг через этот файл, чтобы оценить его соседей, из-за чего упускается часть информации.

В плотно связанных подкорпусах, таких как \texttt{realistic},
типизированный подход обладает большей точностью контекста, но
на разреженных графах доминирует обычный обход, видимо,
потому что он позволяет делать дополнительные шаги через нерелевантные узлы,
что иногда позволяет добраться до необходимой информации.

При этом нельзя сделать однозначный вывод о влиянии учета типов ссылок на корректность ответов:
значения никак не коррелируют.

\subsubsection{Анализ влияния графа}

Результаты в таблице~\ref{tab:exp3_vs_exp5}
уже показали устойчивость и эффективность сгенерированных ссылок.
Данный эксперимент нацелен подтвердить,
что этот результат сохраняется при использовании
более мощных и строгих LLM как для генерации ответа,
так и для оценки.

\begin{table}[H]
    \caption{Сравнение качества работы GraphRAG между размеченным графом и сгенерированным.}
    \label{tab:exp3_vs_exp5}
    \centering
    \begin{tabular}{llccc|ccc}
        \hline
        \multirow{2}{*}{\textbf{Датасет}}   & \multirow{2}{*}{\textbf{Алгоритм}} & \multicolumn{3}{c|}{\textbf{Разметка}} & \multicolumn{3}{c}{\textbf{Сгенерированные}}                                                                               \\
        \cline{3-8}
                                            &                                    & \textbf{CP}                            & \textbf{CR}                                  & \textbf{AC} & \textbf{CP}        & \textbf{CR}        & \textbf{AC}         \\
        \hline
        \multirow{2}{*}{\texttt{hub}}       & Base                               & 0.660                                  & 0.875                                        & 0.698       & 0.689 $\uparrow$   & 0.833 $\downarrow$ & 0.682 $\downarrow$  \\
                                            & Typed                              & 0.707                                  & 0.833                                        & 0.709       & 0.707 $=$          & 0.833 $=$          & 0.716 $\uparrow$    \\
        \hline
        \multirow{2}{*}{\texttt{personal}}  & Base                               & 0.750                                  & 1.000                                        & 0.796       & 0.729 $\downarrow$ & 0.958 $\downarrow$ & 0.762 $\downarrow$  \\
                                            & Typed                              & 0.729                                  & 0.917                                        & 0.770       & 0.729 $=$          & 0.917 $=$          & 0.770 $=$           \\
        \hline
        \multirow{2}{*}{\texttt{realistic}} & Base                               & 0.426                                  & 1.000                                        & 0.738       & 0.426 $=$          & 1.000 $=$          & 0.738 $=$           \\
                                            & Typed                              & 0.491                                  & 0.954                                        & 0.792       & 0.476 $\downarrow$ & 0.977 $\uparrow$   & 0.774 $\downarrow$  \\
        \hline
        \multirow{2}{*}{\texttt{sequence}}  & Base                               & 0.434                                  & 0.765                                        & 0.734       & 0.434 $=$          & 0.765 $=$          & 0.734 $=$           \\
                                            & Typed                              & 0.429                                  & 0.765                                        & 0.706       & 0.429 $=$          & 0.765 $=$          & 0.706 $=$           \\
        \hline
        \multirow{2}{*}{\texttt{versus}}    & Base                               & 0.519                                  & 0.655                                        & 0.623       & 0.519 $=$          & 0.655 $=$          & 0.623 $=$           \\
                                            & Typed                              & 0.547                                  & 0.568                                        & 0.576       & 0.547 $=$          & 0.568 $=$          & 0.576 $=$           \\
        \hline
        \multirow{2}{*}{\textbf{Среднее}}   & Base                               & 0.558                                  & 0.859                                        & 0.718       & 0.559 $\uparrow$   & 0.842 $\downarrow$ & 0.708 $\downarrow$  \\
                                            & Typed                              & 0.581                                  & 0.807                                        & 0.711       & 0.578 $\downarrow$ & 0.812 $\uparrow$   & 0.708  $\downarrow$ \\
        \hline
    \end{tabular}
\end{table}


На основе полученных результатов (таблица~\ref{tab:exp3_vs_exp5})
можно сделать вывод, что использование сгенерированных связей
практически не влияет на качество работы GraphRAG.

Падение метрики AC при переходе на автоматически построенный граф
в среднем не превышает 0.01
(с $0.718$ до $0.708$ для Base-обхода и с $0.711$ до $0.708$ для Typed-обхода).
Самое большое падение AC наблюдается в случае базового алгоритма
в образце \texttt{personal} и составляет 0.034 (с 0.796 до 0.762).

Изменения значений Context Precision и Context Recall
незначительны, что говорит о сопоставимом качестве связей
с ручной разметкой.

\subsection{Выводы и ограничения исследования}
\label{sec:conclusions}

Опыты над алгоритмом построения графа определили оптимальную конфигурацию
извлечения контекста: переход к рекурсивному алгоритму вместе с использованием
Jina как модели переранжирования позволил достичь $F_1 = 0.498$,
что в комбинации с результатами стратегии \texttt{Strict} дает
оптимальный баланс между точностью и полнотой извлеченных фрагментов.

Эксперименты над LLM-классификацией связи доказали критичность использования
few-shot промпта с четкими инструкциями, а также важность современной архитектуры языковой модели:
Gemma-4-31B показала результат Overall $F_1 = 0.669$,
превзойдя более крупную Qwen-3-235B ($F_1 = 0.445$).
Также важным оказалось объединение семантически схожих классов ссылок
для упрощения классификации языковой моделью.

Оценка предложенных алгоритмов GraphRAG через RAGAS
показала их эффективность:
по результатам видно явное превосходство над ObsidianRAG
и сопоставимое с Neo4j качество работы.
Разработанная логика сбора контекста повысила CR по всех подкорпусах вплоть до 1.0.
Сравнение типизированного обхода с обычным показало гибкость системы:
алгоритм может точно извлекать достаточный объем информации в разных сценариях.
Типизированный обход лучше всего подходит для работы с
шумными, плотно связанными графами.
Базовый, в свою очередь, демонстрирует отличные результаты на разреженных графах.

Финальное сравнение работы предложенных алгоритмов GraphRAG
на сгенерированных ссылках показало высокое качество пайплайна для построения графов знаний в Obsidian.
Конвейер устойчив к шуму и может справляться с неоднозначными и сложными ситуациями.

Важно отметить, что во всех экспериментах значения метрики Answer Faithfulness
были не ниже 0.88,
что говорит об отсутствии галлюцинаций и использовании языковыми моделями приведенного контекста.

Несмотря на все достижения,
у проведенного исследования есть ряд ограничений:
\begin{itemize}
    \item Размер и однотипность тестовых данных.
          Датасет содержит всего 5 подкорпусов, которые вместе
          насчитывают 63 заметки и 246 ссылок (таблица~\ref{tab:dataset_samples}).
          Их содержимое в основном~--- статьи Википедии.
          Это сильно ограничивает значимость метрик
          и делает невозможным оценку масштабируемости предложенных подходов
          и эксперименты с гиперпараметрами GraphRAG.
    \item Ограниченность онтологии.
          Набор используемых данных содержит всего несколько схожих тематических групп.
          Исходя из этого, невозможно предположить поведение системы на данных из другой области.
    \item Зависимость от качества разметки.
          Валидационных датасет создавался одним человеком,
          что вносит долю субъективности в типы некоторых связей.
\end{itemize}


\clearpage
\section{Заключение}

По результатам проведенного исследования
можно сформулировать следующие выводы:

Во-первый, на текущий момент на рынке отсутствует
оптимальное решение для автоматизированного создания
ГЗ в Obsidian.
Представленные инструменты имеют компромиссы.
Корпоративные графовые базы данных (Neo4j и MemGraph)
требуют дополнительную инфраструктуру,
чрезмерно зависят от LLM и не поддерживают инкрементальные обновления.
Облачные сервисы для ведения заметок имеют риски утечки данных
и привязывают пользователей к экосистеме.
Эти недостатки подтверждают необходимость создания
легковесного и независимого от поставщиков алгоритма,
сохраняющего приватность пользовательских данных.

Во-вторых, анализ готовых решений позволил
спроектировать и реализовать
эффективный, отказоустойчивый конвейер
для автоматизированного построения ГЗ в Obsidian.
Разработанное решение совмещает гибридный поиск
(BM25 и семантический по плотным эмбеддингам BGE-M3 на базе LanceDB)
с фильтрацией через Cross Encoder (Jina)
и типизацию связей с помощью LLM (Gemma-4-31B).
Алгоритм также имеет модель для сохранения состояния конвейера
(\texttt{MetadataManager}),
позволяющий возобновлять работу алгоритма
с того же этапа,
и состояния хранилища (\texttt{VaultManager}),
проверяющий заметки на наличие изменений
посредством сравнения сохраненного списка файлов
с предыдущего запуска и SHA-256 хеша их содержимого,
и обеспечивающее инкрементальные обновления.
Такой подход позволяет
экономить вычислительные ресурсы,
что крайне важно для индивидуального использования.

В-третьих, для тестирования собственного решения,
в частности для сравнения типизированного GraphRAG и нетипизированным,
был создан модифицированный алгоритм GraphRAG,
основанный на решении ObsidianRAG.
Его особенностями стали учет типов связей при сборе контекста
многошаговый обход графа.

В-четвертых, для оценки системы и сравнения с аналогами
была разработана методология тестирования и
собран валидационный набор данных.
Выбор Википедии как основного источника данных для заметок
позволил соблюсти строгие ограничения касательно
содержимого заметок и структуры хранилища.
Собранный датасет состоит из 63 заметок, имеющих 246 связей,
и разделен на 5 сценариев, каждый из которых
проверяет работоспособность алгоритма построения ГЗ
в одном из сложных случаев,
таких как омонимия, пересекающиеся темы, многошаговая хронология и
недостаток контекста.

В-пятых, был проведен ряд сравнений и экспериментов по улучшения
полученного решения.
Полученные результаты подтвердили высокую эффективность и
точность работы предложенных алгоритмов.
Выявлено, что оптимальной стратегией извлечения контекста
является гибридный поиск: она обеспечивает наилучший баланс,
достигая значения F1-меры 0.680.
Для типизации связей лучше всего использовать
few-shot промпт с детально описанными допустимыми типами и задачей.
Такой подход позволил получить значение взвешенной F1-меры
равным 0.669 (и 0.802 без учета типизации).
Оценка собственного алгоритма GraphRAG
доказала целесообразность улучшений, продемонстрировав
значительный рост метрики Context Recall, получив 0.96 в среднем
и 1.0 на подкорпусах \texttt{hub} и \texttt{personal} в случае
нетипизированного обхода и 0.940 в среднем для типизированного,
при полном отсутствии галлюцинаций
(Answer Faithfulness была не ниже 0.88).
Для сравнения, ObsidianRAG имеет среднее значения CR равное 0.786,
а AC = 0.496.
Важным достижением стало минимальное изменение
полученных значений при переходе от размеченного вручную ГЗ
к сгенерированному: падение CR составило 0.013, AC~--- 0.01
в среднем по всем алгоритмам.
Также было проведено сравнение качества извлечения на сгенерированном
с помощью neo4j Graph Builder.
Это решения имеет наивысшие значения
CR и AC равные 0.977 и 0.602 соответственно.

Разработанный алгоритмов является
практически применимым решением,
устраняющим когнитивное «узкое горлышко»,
заключающееся в ручном связывании многочисленных заметок
в системах управления знаниями.
Конвейер позволяет без участия человека превратить
хранилище в граф знаний, который можно использовать
для эффективного извлечения данных с помощью GraphRAG.
Решение демонстрирует качество работы наравне с
продвинутым решением для корпоративного использования (neo4j).

Несмотря на достижения отличных результатов,
в работе есть ряд ограничений.
Прежде всего это ограниченный объем датасета (63 заметки),
однотипность его содержания и узкая онтология типов связей.
Также недостатком можно считать отсутствие дополнительных
шагов, увеличивающих эффективность извлечения, таких как
создание сообществ.

Дальнейшее развитие решения заключается в
проведении оценки на хранилищах большего размера,
состоящих из сотен или тысяч реальных обезличенных
заметок различных форматов и тематик.
Кроме того, перспективно добавить
генерацию сообществ и автоматическую адаптацию
используемых типов связей (объединение, добавление и удаление типов).
Развитие разработанной системы позволит
создать полностью автономный алгоритм
конструирования ГЗ в Obsidian.


% Представленная выпускная квалификационная работа посвящена
% актуальной проблеме преодоления «когнитивного узкого места»
% в системах управления персональными знаниями (PKM). В условиях
% экспоненциального роста объемов информации классические методы
% ведения заметок и ручного связывания данных становятся неэффективными,
% что ведет к деградации персональных баз знаний. В ходе
% проведенного исследования была успешно достигнута главная
% цель работы: разработан, реализован и всесторонне оценен
% автоматизированный алгоритм семантического связывания
% текстовых данных для экосистемы Obsidian, а также
% усовершенствован конвейер генерации, дополненной поиском по графу знаний (GraphRAG).

% На основании проведенного теоретического и практического
% исследования можно сформулировать следующие основные выводы:

% Во-первых, анализ существующих решений показал, что на
% рынке отсутствует оптимальный инструмент для локальных
% баз знаний, который сочетал бы в себе вычислительную
% эффективность, суверенность данных и устойчивость к шуму.
% Корпоративные графовые базы данных (такие как Neo4j) требуют
% дорогостоящей инфраструктуры и не поддерживают инкрементальные
% обновления без перестройки всего графа, а облачные сервисы создают
% неприемлемые риски утечки приватной информации. Это подтвердило
% острую научно-практическую необходимость в создании собственного
% легковесного алгоритма, независимого от закрытых вендорских экосистем.

% Во-вторых, в рамках практической части был спроектирован и
% реализован отказоустойчивый конвейер построения графа знаний.
% Разработанный алгоритм комбинирует гибридный поиск (
% BM25 и плотные BGE-M3 эмбеддинги на базе LanceDB)
% с фильтрацией посредством Cross-Encoder (Jina) и
% типизацией связей с помощью локальных языковых моделей.
% Ключевым архитектурным достижением стала реализация модуля
% управления состоянием (VaultManager и MetadataManager),
% который сканирует хранилище по хешам SHA-256 и обеспечивает
% инкрементальное обновление графа. Это позволяет возобновлять
% работу алгоритма без повторной траты вычислительных ресурсов,
% что критически важно для потребительского оборудования.

% В-третьих, для интеграции построенного графа в процессы поиска
% ответов был разработан усовершенствованный четырехэтапный алгоритм
% GraphRAG. В отличие от референсного решения (ObsidianRAG),
% предложенный пайплайн включает этапы многошагового обхода соседей
% и строгого отбора контекста (Pruning) с переранжированием. Были
% реализованы два режима обхода: базовый (доминирующий на разреженных
% графах) и типизированный (эффективно фильтрующий шум на плотно
% связанных графах).

% В-четвертых, для объективной валидации системы была разработана
% оригинальная методология тестирования и собран репрезентативный
% набор данных на базе статей Википедии. Отказ от сырых пользовательских
% заметок в пользу энциклопедических данных позволил соблюсти
% строгое правило «одна заметка — одна сущность». Датасет из 63
% заметок и 246 связей, разделенный на пять топологических сценариев
% (звездообразный граф, омонимия, пересекающиеся темы, хронология и
% персональные записи), обеспечил комплексную проверку системы на
% способность разрешать смысловые конфликты и находить многошаговые
% логические связи.

% В-пятых, результаты проведенных экспериментов со строго зафиксированными
% параметрами воспроизводимости эмпирически подтвердили высокую
% эффективность предложенных решений. Выявлено, что гибридная
% стратегия поиска кандидатов (Combined) обеспечивает наилучший
% баланс, достигая значения F1-меры 0.680. На этапе классификации
% связей применение few-shot промптинга позволило модели Gemma-4-31B
% превзойти более ресурсоемкие аналоги, показав взвешенную F1-меру
% 0.669 (и 0.802 без учета типизации). Оценка конвейера GraphRAG
% продемонстрировала способность системы извлекать практически
% исчерпывающий контекст (метрика Context Recall достигла 0.96)
% при полном отсутствии галлюцинаций (Answer Faithfulness стабильно
% превышала 0.88). Важнейшим результатом стало то, что переход от
% использования эталонного (размеченного вручную) графа к графу,
% сгенерированному нашим алгоритмом автоматически, привел к падению
% корректности ответа (Answer Correctness) в среднем менее чем на 0.01.
% Это доказывает высочайшую устойчивость разработанной системы к информационному шуму.

% Общий вывод по работе. Разработанный комплекс алгоритмов представляет
% собой законченное, научно обоснованное и практически применимое решение,
% ликвидирующее «когнитивное узкое место» в системах управления персональными
% знаниями. Алгоритм автоматизирует рутинные процессы семантического связывания,
% сохраняя полный суверенитет пользовательских данных и высокую скорость
% работы на локальных машинах. Использование размеченного графа в связке
% с модифицированным многошаговым пайплайном GraphRAG кратно повышает
% полноту и точность извлечения информации по сравнению с традиционными
% методами гибридного поиска.

% Несмотря на высокие достигнутые показатели, исследование имеет ряд
% объективных ограничений, задающих векторы для дальнейшей научной работы.
% В первую очередь, это ограниченный объем (63 заметки) и однотипность
% валидационного датасета, а также узкая онтология типов связей.

% Перспективы дальнейшего изучения темы лежат в плоскости масштабирования
% системы. Планируется проведение нагрузочного тестирования конвейера на
% хранилищах, содержащих десятки тысяч файлов, с привлечением реальных
% обезличенных пользовательских баз. Кроме того, перспективным направлением
% является расширение и динамическая адаптация словаря онтологий
% (типов связей) под специфику конкретного пользователя, а также внедрение
% механизмов кластеризации LLM-предсказаний для устранения путаницы между
% семантически близкими классами отношений. Развитие предложенных методов
% позволит создать полностью автономного «умного ассистента» для любых локальных
% баз знаний, функционирующего без вмешательства пользователя.

\clearpage
\section{Список литературы}

\textit{На английском языке}
\printbibliography[keyword=eng, heading=none]

\section{Приложение 1. Визуализация датасета}
\label{sec:appendix_plots}

% --- Настройка нумерации для Приложения А ---
\setcounter{figure}{0} % Обнуляем счетчик рисунков
\renewcommand{\thefigure}{1.\arabic{figure}} % Добавляем букву "А." к номеру
% ---------------------------------------------

\begin{figure}[H]
    \centering

    % --- Первая строка ---
    \begin{subfigure}{0.48\textwidth}
        \centering
        \includegraphics[width=\linewidth]{images/light/hub.png}
        \caption{Подкорпус Hub-and-Spoke}
        \label{fig:app_hub}
    \end{subfigure}\hfill
    \begin{subfigure}{0.48\textwidth}
        \centering
        \includegraphics[width=\linewidth]{images/light/realistic.png}
        \caption{Подкорпус Realistic}
        \label{fig:app_realistic}
    \end{subfigure}

    \vspace{0.5em}

    % --- Вторая строка ---
    \begin{subfigure}{0.48\textwidth}
        \centering
        \includegraphics[width=\linewidth]{images/light/versus.png}
        \caption{Подкорпус Versus}
        \label{fig:app_versus}
    \end{subfigure}\hfill
    \begin{subfigure}{0.48\textwidth}
        \centering
        \includegraphics[width=\linewidth]{images/light/sequence.png}
        \caption{Подкорпус Sequence}
        \label{fig:app_sequence}
    \end{subfigure}

    \vspace{0.5em}

    % --- Третья строка (по центру) ---
    \begin{subfigure}{0.48\textwidth}
        \centering
        \includegraphics[width=\linewidth]{images/light/personal.png}
        \caption{Подкорпус Personal}
        \label{fig:app_personal}
    \end{subfigure}

    \caption{Визуальные представления графов знаний подкорпусов датасета}
    \label{fig:all_dataset_graphs}
\end{figure}




\end{document}
